\section{Proofs of the results on the general class}
\label{app:proofs}

This appendix proves the three main theorems.

\subsection{Proof of \Cref{thm:persistence} (persistence on the basin)}

\paragraph{Step 1: Threshold characterization of the basin.}
Define the one-sided sanction-preference thresholds
$\bar x_T,\bar x_S$ as the smallest values such that the following two
monotonicity conditions hold for all states at and above the
respective threshold ($x\in[\bar x_S,1]$ for the sender condition,
$x\in[\bar x_T,1]$ for the target condition):
\begin{itemize}
\item (Sender, ordinal) after substituting the target policy induced by
the value sets, the family
$\widetilde J_{\Vt,\Vs}(s,x)$ satisfies interval dominance in
$(s,x)$ on $[0,1]\times\basin$; single crossing in $(s,x)$ is a
sufficient discharger.
\item (Target combined slope condition) $\Ut(x,\alpha(x),\cdot)+\beta_T\Vt(\F(x,\cdot))$
is non-decreasing in $s$ on $[\max(x,\bar x_T),1]$.
\end{itemize}
The basin is $\basin=[\max(\bar x_T,\bar x_S),1]$. By
\Cref{ass:sender,ass:target}, both monotonicity conditions hold for
every $x\in\basin$.

\paragraph{Step 2: Forward invariance.}
Fix $x_0\in\basin$. We must show $\sigDynTop_{\Vs}(x_0,\alpha(x_0))
\ge\max(x_0,\bar x)$: the sender's dynamic best response sits at or
above the basin pivot. This is a statement about the \emph{level} of
the sender's upper argmax, and it is important to be precise about
which condition delivers it. It is \emph{not} \Cref{ass:sender}. An
ordinal condition --- Milgrom-Shannon single crossing, or the weaker
interval dominance the certificate actually carries (\Cref{rem:ido})
--- constrains how the upper argmax \emph{moves} as a parameter rises;
it does not say where that argmax \emph{lies} relative to
$\max(x_0,\bar x)$ when the continuation $\Vs$ is an arbitrary,
possibly non-monotone, member of the $K_S$-Lipschitz value set
$\mathcal C_S$ --- a member whose continuation payoff may fall by as
much as
$K_S|\F(x,s_2)-\F(x,s_1)|$ when the sanction rises from $s_1$ to
$s_2$.

What Step 2 uses is the \emph{cardinal transition-adjusted slope condition}
\eqref{eq:F-adjusted-sender-chord}, instantiated \emph{below} the
pivot in the pivot-star form of (P5): at each basin state $x$, the sender's
stage-payoff gain from raising any sanction $s<\max(x,\bar x)$ to the pivot
must exceed $\beta_S K_S|\F(x,\max(x,\bar x))-\F(x,s)|$.
(\Cref{rem:wedge-equivalence} records the same
strict inequality on the complementary range $[\max(x,\bar
x),1]$, where what it discharges is the \emph{ordinal} condition. The
level argument needs the inequality on the sub-pivot range instead,
and the sub-pivot range is where the machine-checked certificate
carries it. The two are distinct conditions, and neither is used in
place of the other.) Under the sub-pivot slope condition the stage gain strictly
exceeds the worst-case Lipschitz absorption of the continuation, so
every sanction \emph{below} the pivot gives a strictly lower value of
$J_S(x_0,\alpha(x_0),\cdot;\Vs)$ than the pivot itself, for every
$\Vs\in\mathcal C_S$, and hence no maximizer lies below it
(\Cref{lem:single-crossing}(b) below). That gives
$\sigDynTop_{\Vs}(x_0,\alpha(x_0))\ge\max(x_0,\bar x)$.

Two disclosures follow, and we state them plainly. First, the level
clause is proved under a cardinal condition that \Cref{ass:sender}
does not supply; the machine-checked certificate carries that
condition --- a stage lower-slope field on the sub-pivot range, a
Lipschitz field for the value set, and the transition-adjusted slope condition --- as explicit
hypotheses, alongside its ordinal field. A reader who wishes to keep
the sender-side assumption set purely ordinal must take the level
bound itself as a primitive: an ordinal condition on the parametrized
family fixes the comparative static, not the location of the argmax,
and against a $K_S$-Lipschitz continuation that may decline in $s$ the
location has to be won by a cardinal comparison of the stage gain with
the worst-case continuation loss. Second, the cardinal condition does
\emph{not} displace the ordinal condition: the certificate carries
both, as separate fields, and they discharge different steps --- the
cardinal condition the \emph{level} (this step); the ordinal interval-dominance
condition of \Cref{rem:ido} the \emph{direction}, namely the monotone
upper argmax in the state used in Step 4 below. The comparison at a
fixed target policy (\Cref{thm:dominance}) instead uses monotonicity of the
equilibrium sender value over successor states.

Hence $s^E_0\ge\max(x_0,\bar x)\ge x_0$. By sanction monotonicity and
the diagonal condition,
\[
x^E_1=\F(x_0,s^E_0)\ge\F(x_0,x_0)=x_0\ge\bar x,
\]
so $x^E_1\in\basin$. This is the general-$\F$ argument. Under
\eqref{eq:affine} the displayed inequality becomes
$(1-\delta)x_0+\delta s^E_0\ge x_0$. Induction yields
$x^E_t\in\basin$ for all $t\ge0$.

\paragraph{Step 3: Sanction floor.}
The Step-2 floor argument is stated at an arbitrary basin state, not
only at $x_0$: for every $x\in\basin$ we have
$\sigDynTop_{\Vs}(x,\alpha(x))\ge\max(x,\bar x)$. By Step 2 the
equilibrium path stays in $\basin$, so $s^E_t\ge\max(x^E_t,\bar
x)\ge x^E_t$ for every $t\ge0$. Feeding $s^E_t\ge x^E_t$ into the
transition \eqref{eq:transition} and using the monotonicity (F3)
together with the diagonal normalization (F4) gives
$x^E_{t+1}=\F(x^E_t,s^E_t)\ge\F(x^E_t,x^E_t)=x^E_t$; hence the state
path is non-decreasing and $x^E_t\ge x_0$ for every $t$. Therefore
$s^E_t\ge\max(x^E_t,\bar x)\ge\bar x$, which is conclusion (ii).
The proof actually establishes the stronger inherited-state bound
$s^E_t\ge\max(x_0,\bar x)$. This argument uses only the Step-2 floor
and the transition; in particular it does not presuppose Step 4.

\paragraph{Step 4: Sanction monotonicity.}
We show $s^E_t\le s^E_{t+1}$. The state path is non-decreasing,
$x^E_{t+1}\ge x^E_t$, by Step 3. The sender's dynamic best response
$\sigDynTop_{\Vs}(x,\alpha(x))$ is non-decreasing in $x$ on the basin.
This is the one sender-side step that is purely \emph{ordinal}: it is
the monotone-upper-argmax conclusion applied to the parametric
variation in the state $x$, and in the formalization it is run off the
interval-dominance field the certificate carries (\Cref{rem:ido}), for
which single crossing is sufficient but not necessary. No cardinal
bound enters here. Hence $s^E_{t+1}\ge s^E_t$.

\paragraph{Single-crossing lemma.}
\begin{lemma}[Sender dynamic objective: direction and level]
\label{lem:single-crossing}
Fix a value pair
$(\Vt,\Vs)\in\mathcal C_T\times\mathcal C_S$, with
$\Vs\in\mathrm{Lip}_{K_S}$, and write
$\widetilde J_{\Vt,\Vs}$ for
\eqref{eq:sender-policy-substituted}. Then:
\begin{enumerate}[label=(\alph*)]
\item \emph{(Direction; ordinal.)} Under \Cref{ass:sender},
the greatest maximizer of
$s\mapsto\widetilde J_{\Vt,\Vs}(s,x)$ is non-decreasing in
$x\in\basin$. Single crossing in $(s,x)$ is a sufficient stronger
condition.
\item \emph{(Level; cardinal.)} Write $\hat x=\max(x,\bar x)$. Under the
transition-adjusted slope condition \eqref{eq:F-adjusted-sender-chord}, in its
pivot-star instance $s_1=s<\hat x=s_2$, the upper argmax of
$s\mapsto\widetilde J_{\Vt,\Vs}(s,x)$ on $[0,1]$ lies in $[\hat x,1]$.
\end{enumerate}
Neither clause implies the other. Clause (a) is silent below the
pivot, so it cannot deliver (b); and the stage slope in clause (b)
holds only below the pivot, whereas deriving single crossing on the
whole choice axis $[0,1]$ from the cardinal condition requires the stage
lower slope on all of $[0,1]$, not merely
below the pivot. The certificate accordingly carries a field for each.
\end{lemma}

\begin{proof}
(a) This is the monotone-greatest-argmax implication of the
interval-dominance condition in \Cref{ass:sender}. Milgrom--Shannon
single crossing in $(s,x)$ upcasts to interval dominance with no side
conditions (\Cref{rem:ido}).

(b) For the level
conclusion --- that the upper argmax lies in $[\max(x,\bar x),1]$ ---
we use the cardinal transition-adjusted slope condition on the sub-pivot range,
\emph{not} clause (a) and \emph{not} any monotonicity
of $\Vs$. For $\Vs\in\mathcal C_S=\mathrm{Lip}_{K_S}$ and any
$s<\hat x$, the continuation increment from $s$ to the pivot is
bounded by its worst-case Lipschitz absorption,
$\beta_S\bigl|\Vs(\F(x,\hat x))-\Vs(\F(x,s))\bigr|\le\beta_S
K_S|\F(x,\hat x)-\F(x,s)|$, while
\eqref{eq:F-adjusted-sender-chord} says that the stage gain from $s$ to
$\hat x$ strictly exceeds this allowance; hence
$\widetilde J_{\Vt,\Vs}(s,x)<\widetilde J_{\Vt,\Vs}(\hat x,x)$ for every
$s<\hat x$ and \emph{for every} $\Vs\in\mathcal C_S$ --- using only the
$K_S$-Lipschitz bound, valid even for non-monotone $\Vs$ --- so no
maximizer lies below $\hat x$; no two sub-pivot sanctions are compared.
(This is why the certificate carries \emph{both}
sender-side conditions rather than one, and why the two clauses above
are stated separately. The ordinal interval-dominance field
(\Cref{rem:ido}) delivers the
comparative-statics direction: the monotone upper argmax of Step 4.
The cardinal
stack --- the stage lower slope on the sub-pivot range, the Lipschitz
constant $K_S$ of the value set, and the actual transition displacement ---
separately pins the \emph{level}. Under the affine law it is sufficient
to verify a lower slope $m_S$ and $\beta_S K_S\delta<m_S$. Neither does the
other's job: the ordinal form does not pin the level against a
general, possibly non-monotone, Lipschitz continuation; and the cardinal condition,
in the sub-pivot form the level argument uses, says nothing above the
pivot. Only when the stage lower slope is strengthened to the whole
choice axis $[0,1]$ does the cardinal condition
additionally discharge the ordinal field.)
\end{proof}

Steps 2 and 3 use neither A1 nor the target condition: conclusions (i) and (ii)
rest on the level condition alone, at any equilibrium with a $K_S$-Lipschitz
sender value at which the sub-pivot slope condition holds
(\path{standardCompliance_persistence_generalF_level_star}), and A1 enters only
in Step 4. This completes the proof of \Cref{thm:persistence}.

\subsection{Proof of \Cref{thm:dominance} (dynamic versus myopic sanction at a fixed target policy)}
\label{app:proof-dominance}

Fix any target policy $\alpha^{\circ}\colon[0,1]\to[0,1]$ and any
state $x$; neither an equilibrium nor the basin enters the argument. We
compare the sender's myopic best response and
its dynamic best response at the \emph{same} target
action $\alpha^{\circ}(x)$, holding all other arguments fixed. The
proof is Milgrom--Shannon monotone comparative statics in the sender
continuation value $\Vs$.

\paragraph{Step 1: The dynamic objective single-crosses the myopic one in $s$.}
Fix $x$ and $a=\alpha^{\circ}(x)$. The sender's forward-looking
objective and its myopic ($\Vs\equiv0$) limit are
\[
J_S(x,a,s;\Vs)=\Us(x,a,s)+\beta_S\Vs(\F(x,s)),
\qquad
J_S(x,a,s;0)=\Us(x,a,s),
\]
so their difference is the continuation term alone,
\[
J_S(x,a,s;\Vs)-J_S(x,a,s;0)=\beta_S\,\Vs(\F(x,s)).
\]
The value $\Vs$ is non-decreasing in the state by hypothesis (Step~4),
and $\F(x,\cdot)$ is non-decreasing by \textup{(F3)};
hence $s\mapsto\Vs(\F(x,s))$ is non-decreasing, and the difference
$\beta_S\Vs(\F(x,s))$ is non-decreasing in $s$. A difference non-decreasing in
$s$ is exactly the increasing-differences / single-crossing condition that
Milgrom--Shannon require \citep{MilgromShannon1994} --- and here it is sourced
by the \emph{monotonicity} of $\Vs$, not by the pointwise level of any value
function. (Pointwise order alone is insufficient: a value pointwise above the
myopic $0$ need not have larger $s$-increments; monotonicity of $\Vs$ is the
property that does.)

\paragraph{Step 2: The upper argmax of the dynamic objective lies weakly above the myopic one.}
By Milgrom--Shannon Theorem~4 in its upper-extremal form (equivalently,
Topkis' monotonicity-of-the-maximum theorem \citep{topkis1998}), when an
objective differs from a comparator by a function non-decreasing in $s$, the
upper element of its argmax is at least the comparator's. Specializing to
$a=\alpha^{\circ}(x)$ with comparator $J_S(\cdot\,;0)$:
\[
\max\arg\max_s J_S(x,\alpha^{\circ}(x),s;0)
\ \le\ \max\arg\max_s J_S(x,\alpha^{\circ}(x),s;\Vs).
\]

\paragraph{Step 3: The myopic limit is $\Vs\equiv0$.}
The myopic best response $\sigM(x,a)$ is by definition the
element of the stage-payoff argmax. Because the continuation term
vanishes when $\Vs\equiv0$, this is also the upper element of the
dynamic-objective argmax evaluated at the zero continuation value.
The dynamic best response $\sigDynTop_{\Vs}(x,a)$ is the
corresponding upper element at the continuation value
$\Vs$. Therefore, at $a=\alpha^{\circ}(x)$,
\[
\begin{aligned}
\sigM(x,\alpha^{\circ}(x))
  &=\max\arg\max_s
    J_S(x,\alpha^{\circ}(x),s;\Vs\equiv0),\\
\sigDynTop_{\Vs}(x,\alpha^{\circ}(x))
  &=\max\arg\max_s
    J_S(x,\alpha^{\circ}(x),s;\Vs).
\end{aligned}
\]

\paragraph{Step 4: Monotonicity of $\Vs$ (explicit, and derived on gain-ball invariant sets).}
The single-crossing in Step~1 rests on $\Vs$ being non-decreasing in the
state --- \emph{not} on any sign normalization of $\Us$.  Accordingly
\Cref{thm:dominance} states it explicitly.  At the certificate-complete
example economy it is a theorem about the geometry of the invariant set: the
sender's invariant set is the zero-based gain ball
$\{V:\lVert V\rVert_\infty\le M,\ \mathrm{Lip}(V)\le K,
\ V(y_2)-V(y_1)\ge g_S(y_2-y_1)\text{ for }y_1\le y_2\}$ with $g_S\ge0$, so
every member is
non-decreasing.  Lean records the two-step implication in
\path{monotone_of_mem_valLipGainBall} and
\path{theorem2_dynamic_dominates_myopic_at_same_alpha_of_mem_valLipGainBall},
with the theorem at the concrete example economy
\path{eeCert_same_alpha_dynamic_dominance_from_gain_carrier}.
For other economies, Banach cone preservation gives an alternative discharge
\emph{if} $\Us(\cdot,\alpha(\cdot),s)$ is state-monotone; this is
\path{V_S_monotone_of_integrand_monotone}.  A bare Lipschitz value set
alone supplies neither route, so A1--A3 by themselves must not be read as a
proof of $\Vs$ monotonicity.

\paragraph{Step 5: Combining.}
Steps~1--4 give
\[
\sigM(x,\alpha^{\circ}(x))\le\sigDynTop_{\Vs}(x,\alpha^{\circ}(x)),
\]
which is \eqref{eq:thm-dominance-same-alpha}, at every state and for every
bounded continuous non-decreasing $\Vs$
(\path{theorem2_dynamic_dominates_myopic_at_same_alpha}); in particular it
holds at every equilibrium whose sender value is non-decreasing.

\paragraph{Every stage maximizer, not only the selected one.}
The comparison does not depend on how ties in the stage problem are broken.
Under the same non-decreasing $\Vs$, every maximizer $m$ of the stage payoff
$s\mapsto\Us(x,a,s)$ on $[0,1]$ satisfies $m\le\sigDynTop_{\Vs}(x,a)$: Step~2
with $m$ in place of the selected myopic best response gives a maximizer of
$J_S$ at or above $m$, and $\sigDynTop_{\Vs}$ is the greatest one
(\path{exists_dynamic_argmax_ge_of_stage_max},
\texttt{stage\_max\_le\_$\sigma$\_dyn\_top}). As rules, the greatest myopic
best response to any target policy lies weakly below the dynamic best
response at every state
(\path{stageTopRule_le_dynamicTopRealizedRule}). Every comparison is against
the dynamic best response, the greatest maximizer of $J_S$. The claim that every stage maximizer
lies below every maximizer of $J_S$ is false under a non-decreasing
value alone: with a constant $\Vs$ the two objectives differ by a constant,
so they have the same maximizers, and a tied stage argmax then has a stage
maximizer above a maximizer of $J_S$. The strict
ordering of all maximizers holds under the conditions of
\Cref{prop:strict-dominance}.

\paragraph{Specialization to the different-$\alpha$ form.}
Taking $\alpha^{\circ}=\alphaM$ on the myopic side and
$\alpha^{\circ}=\alpha$ on the dynamic side, the fixed-policy
inequality \eqref{eq:thm-dominance-same-alpha} yields the
different-$\alpha$ corollary under (a) $\alpha(x)\le\alphaM(x)$ at the
states compared, and (b) pointwise monotonicity of $\sigM(x,\cdot)$
non-increasing in the compliance argument. Condition (b) is a
comparative static and needs the cross sign, not the own sign: it
follows from decreasing differences of $\Us$ in $(s,a)$ --- the
``compliance placates'' condition $\partial^2\Us/\partial s\,\partial a\le0$
of \Cref{sec:model} --- which orders the unique stage maximizer under
\Cref{ass:concavity} and, without uniqueness, orders the greatest stage
maximizer (\path{isMaxOn_max_of_monotone_diff}, an exchange inequality).
A nonpositive own derivative $\partial\Us/\partial a\le0$ does not imply
it: $\Us=-(s-a)^2-2a$ has that sign and the unique maximizer $s=a$,
increasing in compliance. Both conditions hold for the standard
Kaempfer--Lowenberg payoff family, whose cross term carries the required
sign, but are extraneous to the substantive Milgrom--Shannon content of
\Cref{thm:dominance}.

\paragraph{Proof of \Cref{prop:strict-dominance} (strict comparison at a fixed target policy).}
Fix $x$, $a$ and $\Vs$. Write $\Gamma(s)=\Vs(\F(x,s))$ for the
sender's continuation value at the successor state $\F(x,s)$ and
$J(s)=\Us(x,a,s)+\beta_S\Gamma(s)$ for the forward-looking objective. The map
$s\mapsto\Us(x,a,s)$ is continuous, so its maximizers on the compact interval
$[0,1]$ form a nonempty closed set with a greatest element $m$
(\path{exists_greatest_stage_maximizer}). Every stage maximizer therefore lies
weakly below $m$, which is the first conclusion. For the second, let $t$ be a
stage maximizer and $\hat s$ a maximizer of $J$ on $[0,1]$, and suppose
$\hat s\le t$, so that $\hat s\le m$. Stage optimality of $m$ gives
$\Us(x,a,\hat s)\le\Us(x,a,m)$, and the comparison of continuation values gives
$\Gamma(\hat s)\le\Gamma(m)$. The model maintains $\beta_S\in[0,1)$, and $\beta_S\ge0$
preserves the second inequality after discounting. So $J(\hat s)\le J(m)$.
Rearranged, the hypothesis at $s'$ says $J(m)<J(s')$, so $J(\hat s)<J(s')$ at a feasible
$s'$, which contradicts the optimality of $\hat s$. Hence $t<\hat s$
(\path{all_stage_lt_dynamic_of_greatest_stage_improvement_of_wedge_le}). The
argument compares $\Gamma$ only at pairs $(d,m)$ with $d\le m$, and it uses
the discount only through $\beta_S\ge0$. It uses no concavity,
differentiability or uniqueness, and it does not claim that a maximizer of $J$
exists. A non-decreasing $\Gamma$ is one way to supply the comparison
(\path{all_stage_lt_dynamic_of_greatest_stage_improvement_mono}). \qed

\subsection{Proof of \Cref{thm:shock} (temporary shocks and their asymptotic effect)}

Let $z_0=0$ and $z_{t+1}=\F(z_t,1)$ be the full-sanction path from the
worst feasible initial state. The drift
$d_{\F}(x)=\F(x,1)-x$ is continuous by \textup{(F1)} and strictly
positive on $[0,\bar x]$ by full-sanction ascent below the basin, which
\textup{(F5)} implies. Compactness therefore gives
an attained positive minimum
\[
\eta_{\F}(\bar x)
=\min_{x\in[0,\bar x]}d_{\F}(x)>0.
\]
As long as $z_t<\bar x$,
\[
z_{t+1}-z_t=d_{\F}(z_t)\ge\eta_{\F}(\bar x).
\]
Induction gives $z_t\ge t\eta_{\F}(\bar x)$ while all preceding
iterates remain below $\bar x$. Consequently
\[
\widehat N_{\F}(\bar x)
=\left\lceil\frac{\bar x}{\eta_{\F}(\bar x)}\right\rceil
\]
must satisfy $z_{\widehat N_{\F}(\bar x)}\ge\bar x$: otherwise every
earlier iterate would be pre-basin, the induction would apply through
$\widehat N_{\F}(\bar x)$, and the definition of the ceiling would
contradict $z_{\widehat N_{\F}(\bar x)}<\bar x$.

Now start the same temporary shock from any $x_0\in[0,1]$.
Monotonicity of $\F$ in the inherited state implies inductively that
the resulting shocked path lies weakly above $z_t$ at every date.
Thus
$x_{\widehat N_{\F}(\bar x)}
\ge z_{\widehat N_{\F}(\bar x)}\ge\bar x$
uniformly over $x_0$. The shocked path is non-decreasing, so every
longer shock also brings the state into the basin. Conclusions (i)--(iii) of
\Cref{thm:persistence}, which do not read the target condition P4, give the stated
post-shock conclusions from that date
(\path{standardCompliance_min_drift_shock_generalF_from_sender_package_of_preBasinAscent}).

Under the partial-adjustment law \eqref{eq:affine},
$1-x_t=(1-\delta)^t(1-x_0)$. Therefore the
$N_*(\bar x,\delta)$ displayed immediately after \Cref{thm:shock} is a valid
uniform choice,
and is generally sharper than the minimum-drift bound because it uses
the exact affine trajectory.

\begin{proof}[Proof of \Cref{cor:shock-transition-strength}]
Write
$d_i(x)=\F_i(x,1)-x$ and let $x_2\in[0,\bar x]$ attain the minimum
$\eta_{\F_2}(\bar x)=d_2(x_2)$. The hypothesis $\F_1(x,1)\le\F_2(x,1)$ is
equivalent to $d_1(x)\le d_2(x)$ on the pre-basin interval. Hence
\[
\eta_{\F_1}(\bar x)
\le d_1(x_2)
\le d_2(x_2)
=\eta_{\F_2}(\bar x).
\]
If $\F_1(x,1)<\F_2(x,1)$ throughout the interval, the middle
inequality is strict. The first law's ascent makes $\eta_{\F_1}(\bar x)>0$,
and that hypothesis passes the ascent to the second law, so both minima
are positive with no adjustment hypothesis on $\F_2$; since $\bar x\ge0$,
\[
\frac{\bar x}{\eta_{\F_2}(\bar x)}
\le
\frac{\bar x}{\eta_{\F_1}(\bar x)}.
\]
Monotonicity of the natural-number ceiling gives the stated weak
duration comparison. It cannot generally be strengthened because the
ceiling is locally constant.

Under the affine law,
\[
d_\delta(x)
=\bigl((1-\delta)x+\delta\bigr)-x
=\delta(1-x).
\]
This function is decreasing in $x$, so its minimum on
$[0,\bar x]$ is attained at $\bar x$ and equals
$\delta(1-\bar x)$. Substitution gives the displayed
minimum-drift duration and its monotonicity in $\delta$.
\end{proof}

\subsection{The class under the lifting condition: class-membership panel and proofs of
\Cref{prop:placable-family}, \Cref{thm:regime-separation},
\Cref{prop:sender-only}, \Cref{prop:target-only} and \Cref{thm:four-cells}}
\label{app:proofs:placable}

Every formula in this subsection is transcribed from the formalization at
the implementation head, with the source's quantifiers and domains; the
declarations are listed in \Cref{app:formalization}. The main text names
the conditions; this is the complete certificate and the argument. The
formalization states every result of \Cref{sec:four-regimes} on the class of
panel~A0, whose stage payoffs are abstract (\path{Examples/ChordEconomy.lean},
the structures \texttt{ChordEconomy} and \texttt{ChordEconomyCs}). Panels A--G
then write the argument for the quadratic--linear family, whose rates are exact
and whose constants have closed forms; each of the family's rates is an instance of
a one-sided rate of panel~A0, and the formalization proves every step with the
one-sided rates, as listed at the end of A0. The family is an instance of the
class (\path{toChord}, \path{toChordCs}), so every family formula below is a
theorem at the family's members. Panel~H is written on the class itself, as
\Cref{prop:target-only} is stated, with the family's closed forms in its
Step~4; panel~I assembles \Cref{thm:four-cells}.

\subsubsection*{A0. The class}

\textbf{Stage payoffs.} $\Us$ and $\Ut$ are jointly continuous. Each condition
holds for arguments in $[0,1]$ and, where an increment is signed, for ordered
arguments; it is a difference-quotient inequality, so no differentiability is assumed. For
the sender (the formalization's S2--S12):
\begin{enumerate}[label=(S\arabic*)]
\item \emph{state Lipschitz at full non-compliance}:
$|\Us(x_2,0,s)-\Us(x_1,0,s)|\le w_{\mathrm{hi}}|x_2-x_1|$;
\item \emph{state lower slope at every compliance}:
$\Us(x_2,a,s)-\Us(x_1,a,s)\ge w_{\mathrm{lo}}(x_2-x_1)$ for $x_1\le x_2$, with
$w_{\mathrm{lo}}>0$;
\item \emph{strong concavity in the sanction}, at every real compliance:
$s\mapsto\Us(x,a,s)$ is $(-q)$-semiconcave on $[0,1]$; $q>0$ is not a separate
condition but follows from the concavity condition below;
\item \emph{curvature ceiling}: $s\mapsto\Us(x,a,s)$ is $\bar Q$-semiconvex on
$[0,1]$;
\item \emph{joint semiconcavity at full non-compliance} in (state, sanction), with
moduli $(\sigma_x,r_{\mathrm{hi}},q)$ in the form of (T3), $\sigma_x\ge0$;
\item \emph{state semiconvexity at full non-compliance}, with modulus $\tau_x\ge0$;
\item \emph{increasing differences at full non-compliance}:
$(\Us(x',0,s')-\Us(x',0,s))-(\Us(x,0,s')-\Us(x,0,s))\ge r_{\mathrm{lo}}(x'-x)(s'-s)$
for $x\le x'$, $s\le s'$, with $r_{\mathrm{lo}}\ge0$;
\item \emph{the cross rate is at most $r_{\mathrm{hi}}$}, in absolute value and
at every compliance;
\item \emph{substitution}:
$(\Us(x,a',s')-\Us(x,a',s))-(\Us(x,a,s')-\Us(x,a,s))\le-k_{\mathrm{lo}}(a'-a)(s'-s)$
for $a\le a'$, $s\le s'$, with $k_{\mathrm{lo}}>0$, and, where a result reads
it, the \emph{upper substitution rate} $k_{\mathrm{hi}}$: the same cross difference is at most
$k_{\mathrm{hi}}(a'-a)|s'-s|$ in absolute value;
\item \emph{upper bound on the marginal return}: $\Us(x,0,t)-\Us(x,0,s)\ge(s-t)\bigl(q(s+t)/2-u_{\mathrm{hi}}\bigr)$
for $t\le s$;
\item \emph{lower bound on the marginal return}: $\Us(x,0,s)-\Us(x,0,t)\ge(s-t)\bigl(u_{\mathrm{lo}}-\bar Q(s+t)/2\bigr)$
for $t\le s$, with $u_{\mathrm{lo}}>0$.
\end{enumerate}
For the target (the formalization's T2--T7):
\begin{enumerate}[label=(P\arabic*)]
\item \emph{compliance cost}:
$-\ell_{\mathrm{hi}}(a_2-a_1)\le\Ut(x,a_2,s)-\Ut(x,a_1,s)\le-\ell_{\mathrm{lo}}(a_2-a_1)$
for $a_1\le a_2$, with $\ell_{\mathrm{lo}}>0$;
\item \emph{sanction pain}:
$p_{\mathrm{lo}}(s_2-s_1)\le\Ut(x,a,s_1)-\Ut(x,a,s_2)\le c\,(s_2-s_1)$ for
$s_1\le s_2$, with $p_{\mathrm{lo}}>0$;
\item \emph{state slope, indexed by the sanction}:
$(v_{\mathrm{lo}}+n_{\mathrm{lo}}s)(x_2-x_1)\le\Ut(x_2,a,s)-\Ut(x_1,a,s)$ for
$x_1\le x_2$, with $v_{\mathrm{lo}}>0$ and $n_{\mathrm{lo}}\ge0$, and
$|\Ut(x_2,a,s)-\Ut(x_1,a,s)|\le(v_{\mathrm{hi}}+n_{\mathrm{hi}}s)|x_2-x_1|$. These
are envelopes on the value of entrenchment at each sanction, not cross rates
(on the family $n_{\mathrm{lo}}=n_{\mathrm{hi}}=n$ is the target's state--sanction
cross term, which softens the pain); the quartic member of \Cref{app:witness} has
$n_{\mathrm{lo}}>n_{\mathrm{hi}}$.
\end{enumerate}
The formalization splits the class. (S8), the upper bound in (P1), the lower
bound in (P2) and the upper envelope in (P3) form a relief extension
(\path{ChordEconomyRelief}, round BC), read by the relief-side results, the
Lipschitz estimates and the sharp sufficient condition for non-compliance; \Cref{prop:relief-participation},
\Cref{prop:branch-invariance}, \Cref{prop:certified-equilibria} and
\Cref{prop:monotone-uniqueness} hold on the base class
(\path{ChordEconomy}), and the enriched subclass extends the relief extension.
Since round BJ the base class is slimmer in three further ways. The two-sided
curvature data --- (S6), and the transition's state-section semiconvexity
$M_{\mathrm{cx}}$ and bound $M_{\mathrm{mix}}$ on mixed differences of panel~A --- form a
two-sided extension (\path{ChordEconomyTwoSided}), which the relief extension
extends; only the two-sided invariant set (the modulus $\sigma_C$ below, and through it
the route through the increasing-differences condition to a monotone response
and $L_{\mathrm{resp}}$) and a
route to a monotone rule from bounds on the target's policy
(\path{Examples/PlacableSenderPolicyClass.lean}) read them. The upper
substitution rate of (S9) is a condition of the base results that read it
(\path{PlacationCeiling}), the sanction floor of
\Cref{prop:certified-equilibria} among them, and a field of the relief
extension. And $q>0$ is a
theorem: the concavity condition puts $q$ above
$\beta_S(\sigma_Vd_{\mathrm{hi}}^2+K_SM_s)$, which is non-negative because the
Riccati closure gives $\sigma_V\ge0$. Every member the formalization exhibits
lies in the two-sided extension and has an upper substitution rate, so the
further generality these changes give the base class is not yet illustrated
by an example.
The myopic best response $\sigma(x,a)$ ($\sigM(x,a)$ in the main text), a field
of the game (\Cref{sec:model}), is the
maximizer of $\Us(x,a,\cdot)$ on $[0,1]$, unique by (S3); on the family it is the
clipped vertex, and in general it has no closed form. The class also requires it
to vary continuously with $(x,a)$; since round BJ this is a condition of the
class rather than of the game (\texttt{ChordEconomy.$\sigma$\_continuous}), kept
because the class's results on the four cases compare with the myopic sender
(\Cref{app:moved-ad}). \emph{Compliance
neutrality}, $\Us(x,a,0)=\Us(x,0,0)$, is a hypothesis of the relief-side results
only (the sender-only case, \Cref{thm:four-cells}, the relief equilibrium);
the all-equilibria floor of \Cref{prop:certified-equilibria} reads its one-sided
pairwise form (a larger concession is never worth more to the sender at zero
sanction) and \Cref{prop:relief-participation} its form against full non-compliance.
All three hold on the family, where $\Us(x,a,0)=wx$.

\textbf{Constants.} With $M_U$ and $M_{T,0}$ the suprema of $|\Us|$ and $|\Ut|$
on the cube (not fields),
\[
K_S=\frac{w_{\mathrm{hi}}}{1-\beta_Sg_{\mathrm{hi}}},\quad
g_S=\frac{w_{\mathrm{lo}}}{1-\beta_Sg_{\mathrm{lo}}},\quad
M_S=\frac{2M_U}{1-\beta_S},\quad
M_T=\frac{2M_{T,0}}{1-\beta_T},\quad
\mu=q-\beta_S\bigl(\sigma_Vd_{\mathrm{hi}}^2+K_SM_s\bigr),
\]
\[
\varrho=r_{\mathrm{hi}}+\beta_S\bigl(\sigma_Vg_{\mathrm{hi}}d_{\mathrm{hi}}+K_SM_m\bigr),\qquad
\sigma_{\mathrm{img}}=\sigma_x+\beta_S\bigl(\sigma_Vg_{\mathrm{hi}}^2+K_SM_x\bigr)+\frac{\varrho^2}{\mu},
\]
\[
A=u_{\mathrm{hi}}+\beta_SK_Sd_{\mathrm{hi}},\qquad
\underline s=\sigma(0,0),\qquad
g_T=\frac{v_{\mathrm{lo}}+n_{\mathrm{lo}}\underline s}{1-\beta_Tg_{\mathrm{lo}}},\qquad
\sigma_C=\frac{\tau_x+\beta_SK_SM_{\mathrm{cx}}}{1-\beta_Sg_{\mathrm{hi}}^2},
\]
\[
M_{\mathrm{tot}}=\bar Q+\beta_S\bigl(\sigma_Cd_{\mathrm{hi}}^2+K_SM_{\mathrm{cs}}\bigr),\qquad
L_{\mathrm{resp}}=r_{\mathrm{hi}}+\beta_S\bigl(\bar\sigma g_{\mathrm{hi}}d_{\mathrm{hi}}
+(K_S+2\bar\sigma)M_{\mathrm{mix}}\bigr),\ \ \bar\sigma=\max(\sigma_V,\sigma_C),
\]
and $L_{\mathrm{loop}}=g_{\mathrm{hi}}+d_{\mathrm{hi}}L_{\mathrm{resp}}/\mu$. The
lower bound $\underline s$ is the myopic sanction at full non-compliance at the lowest state,
below which the response at non-compliance never falls.

\textbf{The certificate.} A member of the class satisfies (curvature) $\mu>0$,
the concavity condition, and (Riccati closure) $\sigma_{\mathrm{img}}\le\sigma_V$,
and nothing else is a condition of membership.
Interiority, $A<q$, and the myopic relief condition,
$\ell_{\mathrm{hi}}\bar Q<k_{\mathrm{lo}}p_{\mathrm{lo}}$, are hypotheses of the
results that read them, not fields (\path{Interior}, \path{StageRelief}; round
BC). Interiority implies myopic interiority $u_{\mathrm{hi}}<q$, that a myopic
sender never sanctions fully (\path{uhi_lt_q}), and the result that the
sender's forward-looking incentive changes behavior, the strict foresight premium above
the stage choice, the myopic equilibrium,
the target-only case and the convergence of the forward-looking path read
only that; the remaining results, among them the interior bound on the
dynamic best response (the cap $A/q<1$), the sender-only case, the relief
side and the four cases, keep interiority. The myopic relief condition is read by
the myopic equilibrium of \Cref{thm:regime-separation} and by the threshold
discount factor $\beta_T^*$, and in the four cases it follows from the forward-looking
relief condition (\path{stageRelief_of_reliefSharp}).
The hypotheses of the propositions that read them are (the lifting condition)
$A\le k_{\mathrm{lo}}$ --- which implies the myopic lifting condition
$u_{\mathrm{hi}}<k_{\mathrm{lo}}$ (\path{uhi_lt_k}), all that the myopic
equilibrium reads; the sender-only case, the relief side and the four cases
keep the lifting condition ---, the entrenchment condition $c\le\beta_Td_{\mathrm{lo}}g_T$, the
one-period ranking $\ell_{\mathrm{hi}}<\underline s\,p_{\mathrm{lo}}$, under which
full compliance beats full non-compliance in the one-period problem, the increasing-differences condition (RM)
with $r_{\mathrm{lo}}$ on its right, and the forward-looking relief condition (RL)
$\ell_{\mathrm{hi}}M_{\mathrm{tot}}<k_{\mathrm{lo}}p_{\mathrm{lo}}$. The concavity
condition $\mu>0$ and the closure admit some $\sigma_V\ge0$ exactly under the discriminant
condition of panel~A with $N=B_0Q+C_0^2$,
$L=(1-\beta_Sg_{\mathrm{hi}}^2)Q+B_0D-2C_0\beta_Sg_{\mathrm{hi}}d_{\mathrm{hi}}$,
$B_0=\sigma_x+\beta_SK_SM_x$ and $C_0=r_{\mathrm{hi}}+\beta_SK_SM_m$
(\path{chSigV_feasible_iff}).

\textbf{Replacements.} Panels A--G use the family's exact rates. On the class
the stage vertex $(u_0+rx)/q$ is replaced by $\sigma(x,0)$ and $u_0/q$
by $\underline s$; the linear bound $u_0-ka+rx\le u_0+r$ on the marginal return by (S10); the
stage state rate $rs+w\in[w,r+w]$ by (S1) and (S2); the decreasing differences
$-ksa$ by (S9); the exact cross rate $r$ by (S7) from below and (S8) from above;
the stage curvature $q$ by (S3) from below and (S4) from above; the costs $\ell a$,
the pain $c-nx$ and the target's state rate $v+ns$ by (P1)--(P3); and
$\Us(x,a,0)=wx$ by compliance neutrality. The sharp rate of sanction reduction uses the
argmax shift under a nonlinear tilt (\path{argmax_shift_ge_tilt_of_semiconvex_interior})
in place of the linear shift, the foresight premium uses the shifted-argmax bound under
$\bar Q$-semiconvexity in place of the exact quadratic, and the zero level is
the least zero of an antitone family (\path{Foundations/ZeroLevel.lean}). The
family's closed forms that have no one-sided counterpart stay on the family:
the target-only steady state $u_0/(q-r)$ and its rate, the comparisons across
economies in $r$, the interest-group microfoundations, and the coefficient box.

\subsubsection*{A. Primitives}

\textbf{Stage payoffs.} With state $x$, compliance $a$ and sanction $s$ in
$[0,1]$,
\[
\Us(x,a,s)=-\tfrac q2 s^2+s(u_0-ka+rx)+wx,
\qquad
\Ut(x,a,s)=-\ell a+s(nx-c)+vx .
\]
Coefficient signs supplied: $u_0>0$, $r\ge0$, $w>0$, $\ell>0$, $n\ge0$, $k>0$,
$v>0$, and $c>n$ (hence $c>0$): the sanction hurts the target at every state.
The lifting condition implies $k>0$, the entrenchment condition with interiority implies $v>0$, and
the myopic relief condition implies $c>n$; since these three are hypotheses of the
results that read them, not conditions of membership, the signs are supplied
(see C). The sign $q>0$ is \emph{derived}, from the concavity condition (see C).
Discount factors
$\beta_S,\beta_T\in(0,1)$.

\textbf{Transition.} $\F$ satisfies \textup{(F1)--(F4)} on the square and,
in addition, five slope and curvature bounds with nonnegative moduli, each stated for
arguments in $[0,1]$; a sixth, \textup{(T6)}, is carried by the enriched
subclass on which the sender-only case, \Cref{thm:four-cells}, the theorem
that both forward-looking incentives change behavior without the one-period
ranking (\path{both_channels_active_of_relief}, after
\Cref{thm:four-cells}) and the relief and no-non-compliance sides of
\Cref{prop:bracket} are stated:
\begin{enumerate}[label=(T\arabic*)]
\item \emph{State slope bound} $0\le g_{\mathrm{lo}}\le g_{\mathrm{hi}}\le1$:
for every $s$ and every $y_1\le y_2$,
$g_{\mathrm{lo}}(y_2-y_1)\le\F(y_2,s)-\F(y_1,s)\le g_{\mathrm{hi}}(y_2-y_1)$.
\item \emph{Sanction slope bound} $0<d_{\mathrm{lo}}\le d_{\mathrm{hi}}$: for
every $x$ and every $s_1\le s_2$,
$d_{\mathrm{lo}}(s_2-s_1)\le\F(x,s_2)-\F(x,s_1)\le d_{\mathrm{hi}}(s_2-s_1)$.
\item \emph{Joint semiconcavity} with moduli $M_x,M_m,M_s\ge0$: for
all $x_1,x_2,s_1,s_2$ and $t\in[0,1]$,
\begin{align*}
&t\F(x_1,s_1)+(1-t)\F(x_2,s_2)-\F\bigl(tx_1+(1-t)x_2,\;ts_1+(1-t)s_2\bigr)\\
&\qquad\le\frac{t(1-t)}2\Bigl(M_x(x_1-x_2)^2+2M_m|x_1-x_2|\,|s_1-s_2|+M_s(s_1-s_2)^2\Bigr).
\end{align*}
In the source this is the predicate \texttt{JointSemiconcaveOn} with
third modulus $-M_s$; the literal content is the display, which bounds how
far $\F$ lies \emph{below} its chords. Nothing about concavity of $\F$ is
assumed; the affine law has all three moduli zero.
\item \emph{State-section semiconvexity} with modulus $M_{\mathrm{cx}}\ge0$:
for every $s$, every $x_1,x_2$ and $t\in[0,1]$,
$\F(tx_1+(1-t)x_2,s)-\bigl(t\F(x_1,s)+(1-t)\F(x_2,s)\bigr)
\le\tfrac{M_{\mathrm{cx}}}2\,t(1-t)(x_1-x_2)^2$,
which bounds how far a state section lies \emph{above} its chords; zero for
convex sections and for the affine law.
\item \emph{Bound on mixed differences} with modulus $M_{\mathrm{mix}}\ge0$: for all
$x_1,x_2,s_1,s_2$,
$\bigl|(\F(x_2,s_2)-\F(x_1,s_2))-(\F(x_2,s_1)-\F(x_1,s_1))\bigr|
\le M_{\mathrm{mix}}|x_2-x_1|\,|s_2-s_1|$;
zero for the affine law and for every additively separable transition.
\item \emph{Sanction-section semiconvexity} with modulus $M_{\mathrm{cs}}\ge0$:
for every $x$, every $s_1,s_2$ and $t\in[0,1]$,
$\F(x,ts_1+(1-t)s_2)-\bigl(t\F(x,s_1)+(1-t)\F(x,s_2)\bigr)
\le\tfrac{M_{\mathrm{cs}}}2\,t(1-t)(s_1-s_2)^2$; zero for the affine law and
for convex sanction sections, $1/5$ at the non-affine member. It is a field
of the enriched subclass (\texttt{ChordEconomyCs}; \texttt{PlacableEconomyFcs}
on the family), not of the class:
it is consumed by the sender-only case, by \Cref{thm:four-cells}, where it
bounds the curvature of
the sender's objective from below, by the theorem after \Cref{thm:four-cells}
that both forward-looking incentives change behavior without the one-period
ranking, and by the relief and no-non-compliance sides of
\Cref{prop:bracket}, through $M_{\mathrm{tot}}$. \Cref{prop:placable-family},
\Cref{thm:regime-separation} and \Cref{prop:target-only} are stated
without it.
\end{enumerate}
A supplied modulus $\sigma_V\ge0$ completes the data; it is closed by the
Riccati inequality in C and appears on no paper-facing statement as a free
parameter. Its existence is part of membership of the class: a member is one for
which some $\sigma_V$ satisfies (curvature) and (Riccati closure), so the
pair is a restriction on the primitives, not a preference coefficient.
\begin{remark}[Why the modulus is supplied]
\label{rem:curvature-obstructions}
Semiconcavity in the sense of \citet{CannarsaSinestrari2004} is the property
the sender's invariant set keeps, and two machine-checked refutations say why it must
be kept by a supplied modulus rather than derived: the maintained transition
class alone admits members whose closed loop preserves no curvature cone
(\path{maxJoinTransition_closedLoop_not_semiconcave},
\path{minMeetTransition_closedLoop_not_semiconvex}), and the affine extreme
point $\F(x,s)=s$ with a single-peaked sender loss, whose stage argmax is
unique, continuous, monotone and Lipschitz, has a closed loop with two kinks
of opposite sign that kill every semiconcavity and every semiconvexity modulus
at once (\path{clampedPeakSender_closedLoop_not_semiconcave},
\path{clampedPeakSender_closedLoop_not_semiconvex}). The same phenomenon
appears for bilevel programs in Example~3 of \citet{ChenLiSo2025}, where the
lower level's argmax destroys the curvature the upper level needs; the
constrained case is left open there, and the supplied modulus with its
Riccati closure is the answer this family gives.
\end{remark}

The concavity condition and the Riccati closure, two inequalities in the
supplied modulus $\sigma_V$, admit some $\sigma_V\ge0$ exactly when $Q>0$,
$0<L<2Q$ and $L^2\ge4DN$ in the primitives alone, where, in this paragraph
only,
\[
Q=q-\beta_SK_SM_s,\qquad D=\beta_Sd_{\mathrm{hi}}^2,\qquad
N=\beta_SK_SM_x\,Q+(r+\beta_SK_SM_m)^2,
\]
\[
L=(1-\beta_Sg_{\mathrm{hi}}^2)\,Q+\beta_SK_SM_x\,D
-2\,(r+\beta_SK_SM_m)\,\beta_Sg_{\mathrm{hi}}d_{\mathrm{hi}}
\]
(\path{sigV_feasible_iff}, whose source also names the products inlined here;
the vertex $\sigma_V=L/(2D)$ then certifies). The modulus need not be
supplied. With the denominator cleared, closure is the quadratic inequality
$D\sigma^2-L\sigma+N\le0$, whose least root
$\sigma_-=(L-\sqrt{L^2-4DN})/(2D)$ is nonnegative, satisfies curvature and
closure, and is at most every feasible modulus (\path{sigVLeast_curv},
\path{sigVLeast_closure}, \path{sigVLeast_least}); the member re-certified at
$\sigma_-$ has a larger $\mu$, a smaller $\varrho$, a smaller response
Lipschitz constant and a smaller sharp loop constant
(\path{LloopSharp_withLeastModulus_le}), so every quantitative bound below is
tightest there. The members supply the rational $\sigma_V=1$ for display
($\sigma_-=10-7\sqrt2$ at the affine member); a different feasible modulus
gives different constants --- at the vertex $\sigma_V=10$ the sharp loop
constant of the affine member is exactly $1$
(\path{placableAffineVertex_LloopSharp}) --- and every bound uses one
consistently chosen modulus. At the members $(Q,L,N,D)=(3,5/2,1/4,1/8)$ and
$(3,1623/700,25/28,9/50)$.

\subsubsection*{B. Derived constants}
\[
K_S=\frac{r+w}{1-\beta_Sg_{\mathrm{hi}}},\quad
g_S=\frac{w}{1-\beta_Sg_{\mathrm{lo}}},\quad
M_S=\frac{q/2+u_0+k+r+w}{1-\beta_S},\quad
\mu=q-\beta_S\bigl(\sigma_Vd_{\mathrm{hi}}^2+K_SM_s\bigr),
\]
\[
\varrho=r+\beta_S\bigl(\sigma_Vg_{\mathrm{hi}}d_{\mathrm{hi}}+K_SM_m\bigr),\qquad
\sigma_{\mathrm{img}}=\beta_S\bigl(\sigma_Vg_{\mathrm{hi}}^2+K_SM_x\bigr)+\frac{\varrho^2}{\mu},
\]
\[
A=u_0+r+\beta_SK_Sd_{\mathrm{hi}},\qquad
g_T=\frac{v+nu_0/q}{1-\beta_Tg_{\mathrm{lo}}},\qquad
M_T=\frac{c+v}{1-\beta_T},\qquad
\sigma_C=\frac{\beta_SK_SM_{\mathrm{cx}}}{1-\beta_Sg_{\mathrm{hi}}^2},
\]
\[
L_{\mathrm{resp}}=r+\beta_S\bigl(\bar\sigma g_{\mathrm{hi}}d_{\mathrm{hi}}
+(K_S+2\bar\sigma)M_{\mathrm{mix}}\bigr),\ \ \bar\sigma=\max(\sigma_V,\sigma_C),
\qquad
L_{\mathrm{loop}}=g_{\mathrm{hi}}+d_{\mathrm{hi}}\cdot\frac{L_{\mathrm{resp}}}{\mu}.
\]
$K_S$ is the state-Lipschitz modulus of the sender's value (stage rate
$r s+w\le r+w$ compounded through the upper bound on the transition's state
slope); $g_S$ its lower state slope (the stage state rate $w$ at $s=0$, through
the lower bound on that slope); $M_S$ the
sup-norm radius; $\mu$ the strong-concavity modulus of the sender's dynamic
objective in the sanction; $A$ an upper bound on the sender's dynamic marginal
return to the first unit of sanction, at full non-compliance (a bound, not
the derivative of a value that need not be differentiable); $g_T$ the
target's lower value slope --- the direct valuation $v$ plus the cross term
$n$ at $u_0/q$, compounded through the lower bound on the transition's state
slope, and at $n=0$ equal
to $v/(1-\beta_Tg_{\mathrm{lo}})$; when $u_0\le q$, as under interiority,
$u_0/q$ is the lower bound $\underline s$ (\path{toChord_sfloor}), below
which the response at non-compliance never falls on either invariant set.

\subsubsection*{C. The certificate}
A member of the family satisfies
\begin{align*}
&\text{(curvature)}& \beta_S\bigl(\sigma_Vd_{\mathrm{hi}}^2+K_SM_s\bigr)&<q
\quad(\text{i.e. }\mu>0),\\
&\text{(Riccati closure)}& \sigma_{\mathrm{img}}&\le\sigma_V,
\end{align*}
together with the signs $k>0$, $v>0$ and $n<c$. Four further conditions are
hypotheses of the propositions that read them, not conditions of membership,
so that each case of the family is stated on the same structure:
\begin{align*}
&\text{(interiority)}& A&<q,\\
&\text{(myopic relief condition)}& \ell q&<k(c-n),\\
&\text{(the lifting condition)}& A&\le k,\\
&\text{(entrenchment condition)}& c&\le\beta_Td_{\mathrm{lo}}g_T .
\end{align*}
As on the class of panel~A0, interiority and the myopic relief condition are
hypotheses rather than fields since round BJ (the predicates
\path{PlacableEconomyF.Interior} and \path{PlacableEconomyF.StageRelief},
hypotheses \path{hInt} and \path{hSR}); the family keeps only the sign of the
second, $n<c$, which that condition implies (\path{hcn_of_stageRelief}). A result
that reads interiority only through $u_0\le q$, under which $\underline s=u_0/q$,
takes that alone in
its bridge to the class (\path{toChord_sfloor}).
The lifting condition is read by the results in which full compliance ends the
sanction (the myopic and sender-only cases, the zero level of
\Cref{app:proofs:placable}); the entrenchment condition by the results in which the
forward-looking target maintains non-compliance (the target-only and forward-looking cases, the
target's invariant set). In the source they are the predicates
\path{ChordEconomy.Placable} and \path{ChordEconomy.LossOffset} (on the family
\path{PlacableEconomyF.Placable} and \path{PlacableEconomyF.LossOffset}, to which
they reduce through \path{toChord}), and the corresponding hypotheses are named
\path{hP} and \path{hTM} throughout.
Derived from the membership conditions and the slope bounds, with nothing
separately supplied: $\sigma_V\ge0$ (Riccati closure with $\mu>0$), and hence
$q>0$ (the concavity condition puts $q$ above a sum of non-negative terms); under
interiority $q>A>0$, as $A\ge u_0>0$, and so $u_0<q$; under the lifting condition
$k\ge A>u_0+r$;
under the entrenchment condition with interiority $v>c-n$ (the condition cleared of its
denominator, using $g_{\mathrm{lo}}+d_{\mathrm{lo}}\le1$ from the two lower
slope bounds and the diagonal, and $u_0<q$);
$g_{\mathrm{lo}}\le g_{\mathrm{hi}}$ and $d_{\mathrm{lo}}\le d_{\mathrm{hi}}$ (the
slope bounds between $0$ and $1$); $0\le M_{\mathrm{mix}}$ (the mixed bound at the
corners); $d_{\mathrm{lo}}\le1$ (the sanction slope bound at state $0$ between sanctions $0$
and $1$); and feasibility of $\F$ on the square (monotonicity and the
diagonal). In the source six of these (feasibility, the two orderings,
$0\le M_{\mathrm{mix}}$, $q>0$, $\sigma_V\ge0$) are theorems of the
structure, not fields (\path{PlacableEconomyF.derived_fields}, from
\path{mem_Icc_of_mono_of_diag}, \path{StateChordBandAll.rates_le},
\path{sanctionChord_rates_le}, \path{MixedChordBound.nonneg},
\path{plSigV_nonneg_of_closure}, and for $q>0$ the concavity condition,
\path{PlacableEconomyF.hq}, in place of interiority, \path{plq_pos_of_interior},
since round BJ); the signs $k>0$, $v>0$ and $n<c$, which the lifting condition, the entrenchment
condition and the myopic relief condition used to imply, are fields since those became
hypotheses; a member supplies the remaining conditions only. The one-period ranking $\ell q<u_0(c-n)$ --- the endpoint comparison under which
full compliance beats full non-compliance in the one-period problem --- is not a field: it
implies the myopic relief condition through the lifting condition ($u_0\le k$), and it is a
hypothesis of the clause of \Cref{prop:placable-family} that the forward-looking
incentives change behavior
(\path{targetStage_one_gt_zero}, \path{both_channels_active_of_margin}) and of
the dischargers \path{reliefSharp_of_scaled_Mtot_le} and
\path{reliefSharp_of_Mtot_eq_q}; on the enriched subclass \textup{(RL)}
implies the myopic relief condition ($q\le M_{\mathrm{tot}}$) and replaces the
one-period ranking in that clause (\path{both_channels_active_of_relief}). The entrenchment
condition is stated at the refined slope $g_T$; the strong entrenchment condition
$c\le\beta_Td_{\mathrm{lo}}v/(1-\beta_Tg_{\mathrm{lo}})$ implies it
(\path{hTMargin_of_old}) and coincides with it at $n=0$, and the inclusion
is strict off the affine law: at $\F=$ \path{quadDiagF} with
$(q,u_0,k,r,w,\ell,c,n,v)=(3,2,3,\tfrac1{100},\tfrac3{10},\tfrac35,1,\tfrac1{100},\tfrac{1199}{300})$
every field holds, the entrenchment condition reads $1\le1201/1200$, and the
strong entrenchment condition fails, $1199/1200<1$ (\path{placableFloor},
\path{exists_floor_member_not_old}). Equivalently
$v\ge c(1-\beta_Tg_{\mathrm{lo}})/(\beta_Td_{\mathrm{lo}})-nu_0/q$
(\path{lossOffset_iff_v_ge}): a sufficient valuation threshold, not a
behavioral boundary (\path{loss_offset_converse_counterexample}, cited in
\Cref{sec:existence:placable}). The increasing-differences condition
\[
\beta_S\bigl(\sigma_Cg_{\mathrm{hi}}d_{\mathrm{hi}}+(K_S+2\sigma_C)M_{\mathrm{mix}}\bigr)\le r
\tag{RM}
\]
is a hypothesis of the composed-equilibrium and separation theorems, never a
field; at $M_{\mathrm{cx}}=M_{\mathrm{mix}}=0$ it reads $0\le r$. It is used
only to make the sender's response at full non-compliance, computed at the
constructed value $\Vs^\ast$, non-decreasing in the state (Step~4 of panel~E,
\path{responsePath_monotoneOn_of_margin}), and the formalization states the
following results under that condition instead: the composed equilibrium
(\path{equilibriumTop}, \path{both_channels_active_of_response_monotone}), the
separation and the convergence of \Cref{thm:regime-separation}
(\path{myopic_dynamic_equilibrium_separation_generalF_stage_of_response_monotone},
\path{dynamicPath_tendsto_of_response_monotone}), \Cref{prop:target-only}(iv)
(\path{target_only_regime_generalF_of_response_monotone}), case $11$ of
\Cref{thm:four-cells} (\path{four_cell_generalF'_of_response_monotone}),
\Cref{prop:steady-gap}, \Cref{prop:welfare-nonidentification},
\Cref{prop:family-welfare-nonidentification} and \Cref{prop:welfare-interval}; the forms under the condition are their corollaries,
and \Cref{prop:activation-not-hysteresis} is still stated under the condition.
The non-decreasing response is a condition on an endogenous object, the
maximizer of the sender's objective at a value that is itself a fixed point,
so it cannot be checked from the primitives alone; (RM) is a condition on the
primitives and the supplied moduli that implies it, and the main text states
the condition for that reason. The condition
is not empty on the class: at a transition with concave state sections
(\path{quadDiagF} with its arguments swapped, \path{quadDiagFSwap}, not
affine at any rate and not $0$-semiconvex in the state) and
$(q,u_0,k,r,w,\ell,c,n,v)=(3,1,3,\tfrac{9}{278},1,1,4,\tfrac1{100},20)$,
every field holds and (RM) holds with equality, its left side
$\beta_S\sigma_Cg_{\mathrm{hi}}d_{\mathrm{hi}}=9/278=r$ strictly positive
(\path{placableRallyBinding}, \path{exists_binding_rally_member}); the
mixed term is zero there. The condition restricts the supplied moduli,
not the game: \path{placableQuadDiag} re-certified with the state-section
modulus $M_{\mathrm{cx}}=4$ violates (RM), its left side $270/287$ exceeding
$r=1/2$, while its response at non-compliance --- the same function as at the modulus
$0$ --- is monotone (\path{exists_member_response_monotone_not_margin}), so the
non-decreasing response is strictly weaker than (RM) for the supplied moduli.
No economy that fails (RM) under every choice of moduli is exhibited, and no
member of the enriched subclass that fails (RM) is exhibited for the forms of
\Cref{thm:four-cells} or of the welfare results; whether the response at the binding member would stay
monotone under a smaller $r$ is open. The forward-looking relief
condition
\[
\ell\,M_{\mathrm{tot}}<k(c-n),\qquad
M_{\mathrm{tot}}=q+\beta_S\bigl(\sigma_Cd_{\mathrm{hi}}^2+K_SM_{\mathrm{cs}}\bigr),
\tag{RL}
\]
is likewise a theorem hypothesis, of the sender-only case and
\Cref{thm:four-cells} on the enriched subclass. It is the sharp form of the
earlier two-branch relief condition $2\ell\max\{M_{\mathrm{tot}},k\}<k(c-n)$, which
implies it (\path{reliefSharp_of_relief}) and under which the machine-checked
headline rows are still stated; at the two example economies of the subclass (RL) reads
$3<12$ and $45/14<12$, on the whole affine coefficient box of
\Cref{app:witness:openbox} it holds with $M_{\mathrm{tot}}=q$, and wherever
$u_0M_{\mathrm{tot}}\le kq$ --- in particular wherever $M_{\mathrm{tot}}=q$,
by the lifting condition --- it follows from the one-period ranking (a hypothesis) alone, through the identity
$q\,(k(c-n)-\ell M_{\mathrm{tot}})=(kq-u_0M_{\mathrm{tot}})(c-n)+M_{\mathrm{tot}}(u_0(c-n)-\ell q)$
(\path{reliefSharp_of_scaled_Mtot_le}, \path{reliefSharp_of_Mtot_eq_q}).
The condition $u_0M_{\mathrm{tot}}\le kq$ is the largest on the other
coefficients under which the one-period ranking discharges (RL) for every compliance cost
(\path{reliefSharp_uniform_iff}): outside it, the midpoint of
$(k(c-n)/M_{\mathrm{tot}},\,u_0(c-n)/q)$ is a cost that passes the one-period ranking and
fails (RL).
The gap between the two forms is inhabited
off the affine law: at $\F=$ \path{quadDiagF} with
$(q,u_0,k,r,w,\ell,c,n,v)=(3,2,3,\tfrac1{100},\tfrac3{10},\tfrac35,1,\tfrac1{100},5)$
every family condition holds, (RL) reads $6393/3500<297/100$, and the two-branch
form fails (\path{placableSharpRelief},
\path{exists_sharp_relief_member_not_old}); there
$u_0M_{\mathrm{tot}}=2131/350\le9=kq$ while $M_{\mathrm{tot}}>q$, so the
larger discharger applies and the $M_{\mathrm{tot}}=q$ one does not
(\path{placableSharpRelief_relief_disch}).

The four restrictions that mention the cross term $n$ --- its sign, the
one-period ranking, the entrenchment condition and (RL) --- hold together exactly when
$\max\{0,n_{\min}\}\le n<\min\{c-\ell q/u_0,\,c-\ell M_{\mathrm{tot}}/k\}$
with $n_{\min}=(c(1-\beta_Tg_{\mathrm{lo}})/(\beta_Td_{\mathrm{lo}})-v)\,q/u_0$
(\path{rally_softening_interval_iff}; with the family's own condition, the
myopic relief condition, in place of the one-period ranking, exactly when
$\max\{0,n_{\min}\}\le n<c-\ell M_{\mathrm{tot}}/k$,
\path{rally_softening_interval_iff'}): a larger cross term $n$ relaxes the
entrenchment condition and tightens the two accommodation conditions. The interval is one of
sufficient conditions, not the set of economies in which the target maintains non-compliance.

\textbf{Economic interpretation and cost.} Curvature and closure say the sender's
own cost curvature out-curves what the continuation can add, so the
forward-looking sanction problem is strictly concave and its solution
unique; they are regularity, not behavior. Interiority says that even at
full non-compliance and the top of the state space the first unit of sanction is
worth less than its marginal cost at $s=1$, so the sender never sanctions
fully. The lifting condition says the target's compliance can buy the whole of that
marginal return, so full compliance is met with no sanction. The myopic relief condition says
that at the stage, compliance is worth buying at the margin: below the
relief level a unit of it purchases $k/q$ units of sanction reduction worth at least
$c-n$ each, more than its cost $\ell$; the one-period ranking $\ell q<u_0(c-n)$ is the
endpoint form of the same comparison, full compliance against full non-compliance. The entrenchment condition says the
target's per-unit sanction pain is covered by the discounted value of the
entrenchment that sanction builds, at the target's lower value slope --- the
direct valuation $v$ plus the cross term $n$ at $u_0/q$, the lower bound below which
the sanction never falls when $u_0\le q$, as under interiority; this is
a substantive statement about a government that values political
entrenchment more than it dislikes the sanction, and it is what makes the
forward-looking target maintain non-compliance. The sender's positive valuation $w$ of the
state is a political objective of the sanctioning government, not a
national-welfare gain. (RM) says the stage complementarity $r$ between the
state and the sanction covers two forward-looking forces against a monotone
response: saturation of the continuation (bounded by $\sigma_C$) and
cross-effectiveness of the transition (bounded by $M_{\mathrm{mix}}$). It is the
price of transition generality, paid only by transitions with concave state
sections or cross-curvature.

\subsubsection*{D. Invariant sets of value functions}
\[
\begin{aligned}
\mathcal C_S=\{V:\ &\|V\|\le M_S,\ V\text{ is }K_S\text{-Lipschitz},\
V(y_2)-V(y_1)\ge g_S(y_2-y_1)\text{ for }y_1\le y_2,\\
&V\circ\operatorname{clip}\text{ is }\sigma_V\text{-semiconcave on }[0,1]\},
\end{aligned}
\]
$\mathcal C_S^2=\mathcal C_S\cap\{V\circ\operatorname{clip}\text{ is }\sigma_C\text{-semiconvex on }[0,1]\}$,
and
$\mathcal C_T=\{V:\ \|V\|\le M_T,\ V(y_2)-V(y_1)\ge g_T(y_2-y_1)\text{ for }y_1\le y_2\}$.
All three are closed in the sup norm and nonempty (the affine value
$y\mapsto g\,y$ with the right slope lies in each).

\begin{proposition}[Two-sided target value]
\label{prop:target-two-sided}
On the relief class, let $\gamma_T=g_{\mathrm{hi}}+d_{\mathrm{hi}}L_{\mathrm{resp}}/\mu$
be the sharp loop constant and suppose $\beta_T\gamma_T<1$. Then the target value
of every equilibrium whose sender value lies in the two-sided invariant set is
$K_T$-Lipschitz in the state, with $K_T=L_U/(1-\beta_T\gamma_T)$ and
$L_U=\max\{v_{\mathrm{hi}},v_{\mathrm{hi}}+n_{\mathrm{hi}}\}+cL_{\mathrm{resp}}/\mu$
($v+n+cL_{\mathrm{resp}}/\mu$ on the family)
(\path{equilibrium_V_T_lipschitz_of_carrier}; Hinderer 4.2(c) through the
invariant closed set of $K_T$-Lipschitz values, \path{targetFix_lipschitz}).
This covers the constructed equilibrium (\path{VTstar_lipschitz}) and, with no
further condition, the non-compliance-profile value (\path{VTdef_lipschitz}). For $s_1\le s_2$
the target's discounted continuation gain from the higher sanction is at most
$\beta_TK_Td_{\mathrm{hi}}(s_2-s_1)$, which needs only $\beta_T\gamma_T<1$
(\path{VTdef_continuation_upper}); with a nondecreasing response it is at least
$\beta_Tg\,d_{\mathrm{lo}}(s_2-s_1)$ at every slope $g\ge0$ satisfying the
slope condition, in particular at $g_T$ under the entrenchment condition
(\path{VTdef_continuation_lower}, \path{VTstar_continuation_bracket}).
$\beta_T\gamma_T<1$ is a hypothesis of the theorems that read it, not a field;
$K_T$ and $L_U$ are closed forms in existing constants.
\end{proposition}

\begin{proposition}[Certified equilibria]
\label{prop:certified-equilibria}
Let $E$ be any equilibrium in $\mathcal E$ of an economy in the class whose target
policy is nonincreasing in the state, on a member at which a larger concession
is never worth more to the sender at zero sanction (automatic on the family).
Then the sender's value has state slope
at least $g_S=w_{\mathrm{lo}}/(1-\beta_Sg_{\mathrm{lo}})$ (\path{V_S_lowerSlope_of_antitone}),
and, if the member also has an upper substitution rate
$k_{\mathrm{hi}}$ (\path{PlacationCeiling}; the family has one, $k$), at every
state the realized sanction is at least
$\min\{(u_{\mathrm{lo}}+r_{\mathrm{lo}}x-k_{\mathrm{hi}}\alpha(x)+\beta_Sg_Sd_{\mathrm{lo}})/\bar Q,\,1\}$;
on the family this is the stage vertex at the policy's compliance plus the
foresight premium $\beta_Sg_Sd_{\mathrm{lo}}/q$
(\path{equilibrium_sanction_ge_wedge_of_antitone}).
The general theorem (\path{senderValue_lowerSlope_of_policy_chord},
\path{Theory/Equilibrium/CertifiedEquilibria.lean}) is stated over the
abstract game: the sender operator at any continuous policy is a contraction
on the whole value space, the lower-slope cone is closed, nonempty and
invariant under a joint state-slope condition, so the unique fixed point lies in it; on
the class, under the pairwise neutrality, that condition reduces to the policy being
nonincreasing in the state (\path{chord_of_antitone}). Nothing is said about equilibria whose compliance
rises with entrenchment, and no equilibrium beyond the constructed ones is
asserted to exist.
\end{proposition}

\subsubsection*{E. Proof of \Cref{prop:placable-family}}

\emph{Step 1: the sender's value.} The sender Bellman operator at full
non-compliance,
\[
(\mathsf T_SV)(x)=\max_{s\in[0,1]}\Us(x,0,s)+\beta_SV(\F(x,s)),
\]
maps $\mathcal C_S$ into itself: the sup bound is the stage bound geometrically
summed; the $K_S$-Lipschitz bound is the stage rate $r+w$ plus
$\beta_SK_Sg_{\mathrm{hi}}$ through (T1); the lower slope $g_S$ is the stage
state rate $w$ plus $\beta_Sg_Sg_{\mathrm{lo}}$ through (T1), the maximum being
taken pointwise in $x$; and the semiconcavity of the image is the Schur
modulus $\sigma_{\mathrm{img}}=\sigma+\varrho^2/\mu$ of a supremum over $s$
of a function that is jointly semiconcave with moduli $(\sigma,\varrho,\mu)$ in
$(x,s)$, which requires the concavity condition $\mu>0$ in the same clause (the
modulus is undefined without it) --- the stage term contributes $(0,r,q)$, the continuation
contributes $\beta_S$ times the composition moduli obtained from (T3) and
the invariant set's $\sigma_V$ --- and the closure inequality makes it at most
$\sigma_V$. The operator is a $\beta_S$-contraction in the sup norm, so it
has a unique fixed point $\Vs^\ast\in\mathcal C_S$ (Banach). Because $\mu>0$,
the objective $s\mapsto\Us(x,a,s)+\beta_S\Vs^\ast(\F(x,s))$ is $\mu$-strongly
concave, so the maximizer $\rho(x,a)$ is unique and, by the maximum theorem,
continuous in $x$; it coincides with the dynamic best response of
\Cref{sec:model}.

\emph{Step 2: the response is interior, zero at full compliance, and above the myopic sanction.}
Cap: for $t\le s$ the stage satisfies
$\Us(x,a,t)-\Us(x,a,s)\ge(s-t)\bigl(q(s+t)/2-(u_0+r)\bigr)$ because the
linear coefficient $u_0-ka+rx$ is at most $u_0+r$, while the continuation is
$\beta_SK_Sd_{\mathrm{hi}}$-Lipschitz in $s$ by (T2); a maximizer of a
$q$-strongly concave function with initial difference quotient at most $L_0=u_0+r$
plus an $L$-Lipschitz perturbation lies at or below $(L_0+L)/q=A/q$, and
$A/q<1$ by interiority. The lifting condition: at $a=1$ and any $s>0$ the objective's
change from $s=0$ is at most $s(u_0+r-k+\beta_SK_Sd_{\mathrm{hi}})-\tfrac q2s^2
=s(A-k)-\tfrac q2s^2<0$, so $\rho(x,1)=0$. Foresight premium: the continuation is
nondecreasing in $s$ ((T2) and $g_S\ge0$), so $\rho(x,a)$ lies weakly above the
stage maximizer; at $a=0$ the stage maximizer is the vertex
$(u_0+rx)/q$, and moving the sanction up from it by
$\varepsilon\le\min\{1-\text{vertex},\ \beta_Sg_Sd_{\mathrm{lo}}/q\}$ costs the
stage $\tfrac q2\varepsilon^2$ and gains the continuation at least
$\beta_Sg_Sd_{\mathrm{lo}}\varepsilon$, which is larger. So the vertex is not
a maximizer of $J_S$, no maximizer of $J_S$ is a stage maximizer (the sender's
forward-looking incentive changes behavior), and $\rho(x,0)>(u_0+rx)/q\ge u_0/q$.

\emph{Step 3: the target maintains non-compliance, and the stage disagrees.} Fix
$\Vt\in\mathcal C_T$ and $\Vs\in\mathcal C_S$. The response is antitone in
$a$ (the stage has decreasing differences $-ksa$ in $(s,a)$, and the argmax
is unique), so raising $a$ from $a_1$ to $a_2$ lowers the sanction by
$\Delta\rho\ge0$. The target's objective
$\Ut(x,a,\rho(x,a))+\beta_T\Vt(\F(x,\rho(x,a)))$ changes by at most
$-\ell(a_2-a_1)+\Delta\rho\,(c-nx)-\beta_Tg_Td_{\mathrm{lo}}\Delta\rho
\le-\ell(a_2-a_1)+\Delta\rho\,(c-\beta_Td_{\mathrm{lo}}g_T)<0$,
using (T2) for the state loss, the lower slope $g_T$ for the value loss, and
the entrenchment condition. So $a=0$ is the unique maximizer at every value pair in
$\mathcal C_T\times\mathcal C_S$. At
the stage, full compliance gives $-\ell+vx$ (the sanction is zero), while
full non-compliance gives $\rho(x,0)(nx-c)+vx\le-(u_0/q)(c-n)+vx<-\ell+vx$ by the
one-period ranking, the hypothesis of that clause of \Cref{prop:placable-family};
so no forward-looking target maximizer is a stage maximizer (the target's
forward-looking incentive changes behavior). Under
\textup{(RL)} on the enriched subclass the same follows without the one-period ranking: a
compliance $a\le\min\{1,\mu\rho(x,0)/(2k)\}$ keeps the response positive, the
sharp rate of sanction reduction gives $\rho(x,0)-\rho(x,a)\ge ka/M_{\mathrm{tot}}$, and the
stage gain is at least $a\,(k(c-nx)/M_{\mathrm{tot}}-\ell)>0$
(\path{target_channel_active_of_relief}).

\emph{Step 4: the monotone response and the composed equilibrium.} Under
(RM), a condition on members of the two-sided extension of panel~A0, to which
the family belongs, on $\mathcal C_S^2$ the sender's objective has increasing
differences in
(sanction; state): the stage cross-difference is exactly $r\,\Delta x\,\Delta s$,
and the continuation's is bounded below by
$-\beta_S\bigl(\sigma_Cg_{\mathrm{hi}}d_{\mathrm{hi}}+(K_S+2\sigma_C)M_{\mathrm{mix}}\bigr)\Delta x\,\Delta s$
through the semiconvexity $\sigma_C$ of $\Vs^\ast\circ\operatorname{clip}$ and
(T1), (T2), (T5). With a unique argmax, increasing differences make
$\rho(\cdot,0)$ nondecreasing on $[0,1]$. This is the only use of (RM): from
here on only that monotonicity is read, and the formalization builds the
composed equilibrium from it (\path{equilibriumTop}; the form under the increasing-differences
condition, \path{equilibriumTop_of_margin}, is its instance). Given that, the target operator at
the response, $(\mathsf T_TV)(x)=\Ut(x,0,\rho(x,0))+\beta_TV(\F(x,\rho(x,0)))$,
maps $\mathcal C_T$ into itself: the sup bound is direct, and the lower slope
$g_T$ is preserved because the stage payoff at non-compliance rises by at least
$v+n\rho\ge v+nu_0/q$ per unit of state (the response never falls below
$u_0/q$, \path{responsePath_ge_sfloor}), and a higher state gives a weakly
higher sanction and hence, through (T1) and the entrenchment condition, a continuation
gain that at least offsets the stage loss $-(c-nx)\Delta\rho$. It is a
$\beta_T$-contraction;
its fixed point $\Vt^\ast$ and $\alpha\equiv0$ complete
$(\Vt^\ast,\Vs^\ast,0)\in\mathcal E$, with both forward-looking incentives changing behavior by Steps 2
and 3. $\Vs^\ast\in\mathcal C_S^2$ because $\sigma_C$ is the fixed point of
the linear closure $\sigma\mapsto\beta_S(\sigma g_{\mathrm{hi}}^2+K_SM_{\mathrm{cx}})$.

\subsubsection*{F. Proof of \Cref{prop:myopic-benchmark} and \Cref{thm:regime-separation}}

\emph{Step 5: the myopic case.} With $\beta_S=\beta_T=0$ the sender's
problem is the stage problem, whose unique maximizer is the clipped vertex
$\max\{0,\min\{1,(u_0-ka+rx)/q\}\}$; it is zero exactly when
$a\ge a^{00}(x):=(u_0+rx)/k$, which lies in $(0,1)$ because $u_0+r<k$. For
$a\ge a^{00}(x)$ the target's payoff is $-\ell a+vx$, strictly decreasing;
for $a<a^{00}(x)$ it is $-\ell a+\tfrac{u_0-ka+rx}q(nx-c)+vx$, whose slope
in $a$ is $-\ell+\tfrac kq(c-nx)>0$ because $\ell q<k(c-n)\le k(c-nx)$, the
myopic relief condition (\path{myopic_target_lt}). So
$a^{00}(x)$ is the unique optimum, the realized sanction is zero, and the
values $wx$ and $vx-\ell a^{00}(x)$ satisfy the two Bellman equations of
\Cref{sec:model} at zero discount, where the dynamic best responses reduce to the
myopic best response because the stage objective is strictly concave in the
sanction. Convergence to zero: (T2) at state $x$ between sanctions $0$ and $x$, with
$\F(x,x)=x$, gives $\F(x,0)\le(1-d_{\mathrm{lo}})x$, hence
$x_t\le(1-d_{\mathrm{lo}})^tx_0$.

\emph{Step 6: the forward-looking paths.} Write $b=u_0/q$ and
$\rho(y)=\rho(y,0)\in(b,A/q]$ by Step 2. The lower bound on the build-up of the state: when $y\le b$, (T2)
between sanctions $y$ and $b$ gives $\F(y,\rho(y))\ge\F(y,b)\ge y+d_{\mathrm{lo}}(b-y)$,
and the comparison recursion $z_{t+1}=z_t+d_{\mathrm{lo}}(b-z_t)$ from $z_0=0$
gives $x_t\ge b\bigl(1-(1-d_{\mathrm{lo}})^t\bigr)$; the affine minorant is
used only at $y\le b$. From a common start the forward-looking state strictly
exceeds the myopic one: at $t=1$,
$\F(x_0,\rho(x_0))-\F(x_0,0)\ge d_{\mathrm{lo}}\rho(x_0)>0$, and
monotonicity of $\F$ in the state propagates the strict gap. None of this
uses monotonicity of $\rho$.

\emph{Step 7: convergence.} The response $\rho$ is continuous without further conditions;
when it is nondecreasing, as Step~4 gives under (RM), the closed loop
$G(y)=\F(y,\rho(y))$ is nondecreasing and continuous on $[0,1]$; its orbit from any $x_0$ is monotone, hence
convergent to a fixed point $x^{11}_\infty$. Since $\F(x^{11}_\infty,\rho(x^{11}_\infty))=x^{11}_\infty=\F(x^{11}_\infty,x^{11}_\infty)$ and
$d_{\mathrm{lo}}>0$, (T2) forces $\rho(x^{11}_\infty)=x^{11}_\infty$, and $\rho$'s range gives
$b<x^{11}_\infty\le A/q$ (\path{dynamicPath_tendsto_of_response_monotone}). The limit
may depend on $x_0$; what is unique is the
\emph{state} at which sanction equals state along that orbit, not the
equilibrium.

\emph{Step 8: the contraction bound (separate).} By strong concavity
$\mu$ and the fact that the objective's dependence on the state is
$L_{\mathrm{resp}}|x-y|$-Lipschitz in the sanction (through $r$, (T1), (T2),
(T5) and $\bar\sigma$), the maximizers satisfy
$|\rho(x)-\rho(y)|\le L_{\mathrm{resp}}|x-y|/\mu$ (the sharp argmax
stability bound, from the strong-concavity inequality at $t\rho(x)+(1-t)\rho(y)$ and $t\downarrow0$;
the midpoint inequality of the first formalization gave the factor four); then
$|G(x)-G(y)|\le g_{\mathrm{hi}}|x-y|+d_{\mathrm{hi}}|\rho(x)-\rho(y)|
\le L_{\mathrm{loop}}|x-y|$. When $L_{\mathrm{loop}}<1$ the steady state is
unique and globally stable (\Cref{prop:activation-not-hysteresis}); this
is not an ingredient of Steps 1--7.

\textbf{Two views of one concept.} Write $E^{00}\in\mathcal E(p^{00})$ for
the myopic case and $E^{11}\in\mathcal E(p)$ for the forward-looking one;
the game is the same, only the discount factors differ. The
comparison at a fixed target policy of \Cref{sec:results} and
this equilibrium comparison, in which both governments re-optimize, are distinct. The two mixed cases are
\Cref{prop:target-only} and \Cref{prop:sender-only}, and
\Cref{thm:four-cells} orders all four; no interaction theorem is claimed
beyond those directions.

\subsubsection*{G. Proof of \Cref{prop:sender-only}}

\emph{Step 1 (the zero-sanction value).} Let
$\mathcal R V(x)=w\,x+\beta_S\,V(\F(x,0))$ on the two-sided sender invariant set of
panel D (sup-norm at most $M_S$, state-Lipschitz $K_S$, lower slope $g_S$,
$\sigma_V$-semiconcave and $\sigma_C$-semiconvex through the clip). $\mathcal
R$ is a $\beta_S$-contraction (\path{relaxOp_contracting}) and maps that
set to itself (\path{relaxOp_mem2}): the norm and Lipschitz legs use the
upper bound on the transition's state slope at $s=0$; the lower-slope leg uses
the lower bound on that slope;
semiconcavity of the composition $V\circ\F(\cdot,0)$ follows from the joint
semiconcavity (T3) restricted to $s_1=s_2=0$, with composed modulus
$\sigma_x+\beta_S(\sigma_Vg_{\mathrm{hi}}^2+K_SM_x)\le\sigma_V$ by the Riccati closure
(\path{chSig_le_SigV}); semiconvexity follows from the state-section
semiconvexity (T4) through the composition lemma of
\path{Foundations/SemiconvexComposition}. Its fixed point $\Vs^{01}$
(\path{VS01}, \path{VS01_fix}) is the value of the path along which the
sanction is zero forever. Nothing about the target has been used: this is why
the sender's value can be written down before the target's policy is known.

\emph{Step 2 (the response and its two curvatures).} At $\Vs^{01}$ the sender's
dynamic objective in the sanction is
$\psi_{x,a}(s)=\Us(x,a,s)+\beta_S\Vs^{01}(\F(x,s))$. It is $\mu$-strongly
concave, as in panel E, so its maximizer $\rho(x,a)$ is unique and, by the
maximum theorem, jointly continuous on $[0,1]\times\mathbb R$
($\rho$\path{01_continuousOn}). It is also $M_{\mathrm{tot}}$-semiconvex in $s$
with
$M_{\mathrm{tot}}=q+\beta_S(\sigma_Cd_{\mathrm{hi}}^2+K_SM_{\mathrm{cs}})$
(\path{senderObj_semiconvex}): the stage term contributes $q$, and the
continuation term is the composition of a $\sigma_C$-semiconvex, $K_S$-Lipschitz
value with a sanction section whose slopes lie in $[d_{\mathrm{lo}},d_{\mathrm{hi}}]$
and which is $M_{\mathrm{cs}}$-semiconvex (T6). This is the only step of this proof
that uses (T6); the same ceiling $M_{\mathrm{tot}}$ enters the relief and
no-non-compliance sides of \Cref{prop:bracket}.

\emph{Step 3 (the rate of sanction reduction).} Since $\Us$ is linear in $a$ with slope $-ks$,
$\psi_{x,a}=\psi_{x,a'}+k(a'-a)\,s$ for $a<a'$. The shift lemma of
\path{Foundations/SemiconvexArgmaxShift} says that if $\psi$ is
$M$-semiconvex on $[0,1]$ with an interior maximizer $s'$, and $s$ maximizes
$\psi+\lambda u$ with $\lambda>0$, then $s-s'\ge\min\{\lambda/M,\,1-s'\}$,
and $s-s'\ge\lambda/M$ outright when $s<1$: at an interior maximizer $z$ of
an $M$-semiconvex function the semiconvexity inequality through $z-\eta$ and $y$,
read with $\eta\to0$ on both sides of $z$, is the quadratic lower support
$\psi(y)\ge\psi(z)-\tfrac M2(y-z)^2$; applying it to $\psi$ at $s'$ and to
$\psi+\lambda u$ at $s$, each evaluated at the other maximizer, and adding
gives $\lambda(s-s')\le M(s-s')^2$, where $s-s'>0$ by the coarse shift
lemma, so that dividing gives $s-s'\ge\lambda/M$
(\path{argmax_shift_ge_of_semiconvex_linear_sharp}; the quadratic
$-\tfrac M2(u-z)^2$ attains the bound, so nothing sharper holds on the
class). Hence, for $a<a'<a^{01}(x)$, where both responses are interior
($\rho(x,a')>0$ below the zero level and $\rho<1$ by interiority),
\[
\rho(x,a)-\rho(x,a')\;\ge\;\frac{k(a'-a)}{M_{\mathrm{tot}}}
\]
($\rho$\path{01_shift_ge_linear}).

\emph{Step 4 (the zero level).} Let $Z(x)=\{a\in[0,1]:\rho(x,a)=0\}$. It is
nonempty by the lifting condition ($\rho(x,1)=0$) and closed by continuity, so
$a^{01}(x)=\min Z(x)$ exists (\path{zeroLevel}, \path{zeroLevel_mem}).
The response is antitone in $a$, so $\rho(x,a)=0$ for $a\ge a^{01}(x)$ and
$\rho(x,a)>0$ for $a<a^{01}(x)$ ($\rho$\path{01_eq_zero_of_ge},
$\rho$\path{01_pos_of_lt}). At the myopic relief level $a^{00}(x)=(u_0+rx)/k$
the stage marginal at $s=0$ vanishes but the continuation marginal does not:
$\Vs^{01}$ has lower slope $g_S>0$ and the sanction section has lower slope bound
$d_{\mathrm{lo}}>0$, so a small positive sanction beats zero
($\rho$\path{01_pos_at_myopicPolicy}). Hence $a^{00}(x)<a^{01}(x)$
(\path{myopicPolicy_lt_zeroLevel}); this is the foresight premium. Letting
$a'\uparrow a^{01}(x)$ in Step 3, with $\rho(x,a')\ge0$, gives for $a<a^{01}(x)$
($\rho$\path{01_ge_of_lt_sharp}; no continuity of the response is needed)
\[
\rho(x,a)\;\ge\;\frac{k\,(a^{01}(x)-a)}{M_{\mathrm{tot}}}.
\]
This is the untruncated per-unit rate of sanction reduction.

\emph{Step 5 (the target's optimum).} With $\beta_T=0$ the target maximizes
$\tau_x(a)=-\ell a+\rho(x,a)(nx-c)+v\,x$ over $[0,1]$
(\path{tObj01}). For $a>a^{01}(x)$, $\rho=0$ and $\tau_x(a)=-\ell a+v\,x<\tau_x(a^{01}(x))$
(\path{tObj01_lt_of_gt}). For $a<a^{01}(x)$, write $\Delta=a^{01}(x)-a\in(0,1]$;
then, by Step~4 and $c-nx\ge c-n>0$,
\[
\tau_x(a^{01}(x))-\tau_x(a)=(c-nx)\rho(x,a)-\ell\Delta
\;\ge\;\Delta\Bigl(\frac{k(c-n)}{M_{\mathrm{tot}}}-\ell\Bigr)>0
\]
by the forward-looking relief condition (RL) $\ell M_{\mathrm{tot}}<k(c-n)$
(\path{tObj01_lt_of_lt'}). So $a^{01}(x)$ is the unique maximizer
(\path{zeroLevel_isMaxOn'}, \path{zeroLevel_argmax_unique'}). (RL) is
a per-unit benefit--cost condition: the sanction pain avoided per unit of
compliance, at the rate of sanction reduction $k/M_{\mathrm{tot}}$ and the lowest price
$c-n$, exceeds the unit compliance cost $\ell$. The two-branch form
$2\ell\max\{M_{\mathrm{tot}},k\}<k(c-n)$ under which the headline rows are
stated implies it (\path{reliefSharp_of_relief}), and those rows are
recovered as corollaries (\path{tObj01_lt_of_lt}, \path{zeroLevel_isMaxOn}).

\emph{Step 6 (continuity of the zero level).} Fix $x_0$ and $\varepsilon>0$.
If $a^{01}(x_0)-\varepsilon<0$ the lower bound $a^{01}(x)>a^{01}(x_0)-\varepsilon$
holds by feasibility; otherwise $\rho(x_0,a^{01}(x_0)-\varepsilon)>0$, and joint
continuity gives
$\rho(x,a^{01}(x_0)-\varepsilon)>0$ for $x$ near $x_0$, so $a^{01}(x)>a^{01}(x_0)-\varepsilon$.
Conversely, if $a^{01}(x)>a^{01}(x_0)+\varepsilon$ then Step 4 gives
$\rho(x,a^{01}(x_0))\ge k\varepsilon/M_{\mathrm{tot}}>0$, while
$\rho(x,a^{01}(x_0))\to\rho(x_0,a^{01}(x_0))=0$; so $a^{01}(x)\le a^{01}(x_0)+\varepsilon$ for
$x$ near $x_0$ (\path{zeroLevel_continuousOn}, from the general continuity of the
least zero under a linear rate of sanction reduction, \path{ZeroLevel.level_continuousOn}). The policy is extended to
$\mathbb R$ through the clip ($\alpha$\path{01}, $\alpha$\path{01_continuous}).

\emph{Step 7 (the two Bellman equations).} Take $V_S=\Vs^{01}$,
$\alpha=a^{01}$ and $V_T(x)=v\,x-\ell a^{01}(x)$ (\path{VT01}). The sender's
equation holds because the maximizer of $\psi_{x,a^{01}(x)}$ is $\rho(x,a^{01}(x))=0$,
so the maximum equals $\Us(x,a^{01}(x),0)+\beta_S\Vs^{01}(\F(x,0))=w\,x+\beta_S\Vs^{01}(\F(x,0))=\Vs^{01}(x)$
by Step 1. The target's equation holds because $\beta_T=0$ and
$\tau_x(a^{01}(x))=-\ell a^{01}(x)+v\,x$ is the maximum by Step 5, which is also the
selector-consistency condition. The tuple is an equilibrium of the sender-only case
(\path{senderOnlyEquilibriumTop}); the equilibrium is not assumed anywhere in
Steps 1--6, which is what Rule 1 requires of it.

\emph{Step 8 (convergence to zero).} The realized sanction is $\rho(x,a^{01}(x))=0$
(\path{senderOnlyRule_eq_zero}), so $x_{t+1}=\F(x_t,0)$. The lower bound on the
transition's sanction slope, between $0$ and $x_t$ and with $\F(x_t,x_t)=x_t$, gives
$x_t-\F(x_t,0)\ge d_{\mathrm{lo}}x_t$, hence $x_t\le(1-d_{\mathrm{lo}})^tx_0\to0$
(\path{senderOnlyPath_le_geom}, \path{senderOnlyPath_tendsto_zero}). \qed

\subsubsection*{H. Proof of \Cref{prop:target-only}}

Clauses (i)--(iii) hold on the class of panel~A0 under the entrenchment condition alone
(\path{target_only_regime_generalF_core}); (iv) adds $u_{\mathrm{hi}}<q$ and
(RM), on the two-sided extension (\path{target_only_regime_generalF}, the
corollary of \path{target_only_regime_generalF_of_response_monotone}); the
closed forms of (iii) are on the family under interiority and the entrenchment condition
(\path{PlacableEconomyF.targetOnlyPath_tendsto},
\path{targetOnlyPath_error_le_geometric}), and the family bundle
\path{target_only_regime_closedForm} adds (RM) for (iv). (T6) is not used. This panel is written
on the class: $\sigma(x,a)$ is the myopic best response and $\underline s=\sigma(0,0)$.

\emph{Step 1 (the sender).} At $\beta_S=0$ the sender maximizes $\Us(x,a,\cdot)$,
which is strictly concave by (S3) (\path{stage_strictConcaveOn}), so
$\sigDynTop_{\Vs}(x,a)=\sigma(x,a)$ for every continuation value
(\path{Theory/Mechanism/ZeroDiscountSelector.lean}). The sanction rule is
therefore $\sigma(x,0)$ (\path{targetOnlyRule_eq_stage}), and at
$\alpha\equiv0$ the sender's equation holds at the stage value
$\Vs^{10}(x)=\Us(x,0,\sigma(x,0))$ (\path{VS10}). None of the following needs a
theorem-specific condition. $\sigma(x,\cdot)$ is antitone, by (S9) and uniqueness
(\path{stage_}$\sigma$\path{_antitoneOn}). $\sigma(\cdot,0)$ is continuous
(panel~A0) and nondecreasing, by (S7) and uniqueness (\path{stage_argmax_mono}),
so it never falls below $\underline s$ (\path{sfloor_le}).

\emph{Step 2 (the target and the equilibrium).} The argument of Step~3 of
panel~E applies with $\sigma(x,\cdot)$ in place of $\rho(x,\cdot)$. Take
$a_1<a_2$ in $[0,1]$, $\sigma_i=\sigma(x,a_i)$ and $\Vt\in\mathcal C_T$. By (P1)
and (P2) the stage term changes by at most
$-\ell_{\mathrm{lo}}(a_2-a_1)+c(\sigma_1-\sigma_2)$, and by (T2) and the slope
$g_T$ the continuation changes by at most
$-\beta_Td_{\mathrm{lo}}g_T(\sigma_1-\sigma_2)$. The objective
$\Ut(x,a,\sigma(x,a))+\beta_T\Vt(\F(x,\sigma(x,a)))$ therefore changes by at most
$-\ell_{\mathrm{lo}}(a_2-a_1)+(\sigma_1-\sigma_2)(c-\beta_Td_{\mathrm{lo}}g_T)<0$,
by the entrenchment condition (\path{targetObj10_strictAntiOn}), so $a=0$ is the unique
maximizer. The operator
$(\mathsf T_T^{10}V)(x)=\Ut(x,0,\sigma(x,0))+\beta_TV(\F(x,\sigma(x,0)))$ is a
$\beta_T$-contraction (\path{targetOp10_contracting}). It maps $\mathcal C_T$
into itself as in Step~4 of panel~E with $\sigma(\cdot,0)$ in place of
$\rho(\cdot,0)$ (\path{targetOp10_mem}). The sup bound is direct. For
$y_1\le y_2$, now write $\sigma_i=\sigma(y_i,0)$; the stage term rises by at least
$(v_{\mathrm{lo}}+n_{\mathrm{lo}}\underline s)(y_2-y_1)-c(\sigma_2-\sigma_1)$
((P3) at $\sigma_1\ge\underline s$, then (P2)), and the continuation rises by at
least
$\beta_Tg_T\bigl(g_{\mathrm{lo}}(y_2-y_1)+d_{\mathrm{lo}}(\sigma_2-\sigma_1)\bigr)$
((T1), (T2)). With $g_T(1-\beta_Tg_{\mathrm{lo}})=v_{\mathrm{lo}}+n_{\mathrm{lo}}\underline s$
and the entrenchment condition the sum is at least $g_T(y_2-y_1)$. Step~4 of panel~E gets
the order $\sigma_1\le\sigma_2$ from (RM); here it comes from (S7). The fixed
point $\Vt^{10}\in\mathcal C_T$ (\path{VT10}), with $\Vs^{10}$ and
$\alpha\equiv0$, is an equilibrium in $\mathcal E$ of the target-only case
(\path{targetOnlyEquilibriumTop}), and at its values every target maximizer is
$0$ (\path{targetOnlyEquilibriumTop_argmax_unique}). This proves (i) and (ii).

\emph{Step 3 (the path, on the class).} The rule is at least $\underline s$, so
Step~6 of panel~F with $b=\underline s$ gives
$x_t\ge\underline s\bigl(1-(1-d_{\mathrm{lo}})^t\bigr)$
(\path{targetOnlyPath_ge_approach}). The rule is nondecreasing and continuous,
so, as in Step~7 of panel~F, the orbit of $y\mapsto\F(y,\sigma(y,0))$ converges
to some $x^{10}_\infty$. (T2) with $d_{\mathrm{lo}}>0$ then forces
$\sigma(x^{10}_\infty,0)=x^{10}_\infty$, and so $x^{10}_\infty\ge\underline s$
(\path{ChordEconomy.targetOnlyPath_tendsto}). No theorem-specific condition enters. The class
result does not say the fixed point is unique, so $x^{10}_\infty$ may depend on $x_0$.

\emph{Step 4 (the family).} Under interiority the myopic best response at non-compliance is
the vertex $(u_0+rx)/q\in[0,1)$ (\path{targetOnlyRule_eq_vertex}). The family's
entrenchment condition is the class's, because $u_0\le q$ gives $\underline s=u_0/q$
(\path{toChord_sfloor}, \path{toChord_lossOffset}). Since $r<q$, the vertex
fixes only $u_0/(q-r)$ (\path{vertex_eq_self_iff}), so $x^{10}_\infty=u_0/(q-r)$
from every start (\path{PlacableEconomyF.targetOnlyPath_tendsto}). For the rate,
the rule is affine about $x^{10}_\infty$ with slope $r/q\in[0,1)$:
$(u_0+rx)/q-x^{10}_\infty=(r/q)(x-x^{10}_\infty)$ (\path{vertex_sub_steady10}). Suppose
$x\le x^{10}_\infty$. Then $R=(u_0+rx)/q$ lies in $[x,x^{10}_\infty]$, monotonicity of $\F$
and the diagonal put $\F(x,R)$ in $[x,R]$, and (T2) gives
$\F(x,R)-x\ge d_{\mathrm{lo}}(R-x)=d_{\mathrm{lo}}(1-r/q)(x^{10}_\infty-x)$. Hence
$|\F(x,R)-x^{10}_\infty|\le\bigl(1-d_{\mathrm{lo}}(1-r/q)\bigr)|x-x^{10}_\infty|$; the argument for
$x\ge x^{10}_\infty$ is symmetric. Iterating gives the centered bound of (iii)
(\path{CenteredContraction.iterate_abs_le},
\path{targetOnlyPath_error_le_geometric}). The bound compares each state with
$x^{10}_\infty$; it does not compare two paths. This proves (iii).

\emph{Step 5 (the forward-looking limit is higher).} Let $u_{\mathrm{hi}}<q$, and
let the response $\rho=\rho(\cdot,0)$ at $\Vs^\ast$ be nondecreasing, which
Step~4 of panel~E gives under (RM). From $x_0$ the forward-looking path
converges to some $x^{11}_\infty$ with $\rho(x^{11}_\infty)=x^{11}_\infty$ (Step~7 of panel~F,
\path{dynamicPath_tendsto_of_response_monotone}). The strict foresight premium of Step~2 of
panel~E holds on the class under $u_{\mathrm{hi}}<q$ alone, with $\bar Q$ in
place of $q$: $\sigma(x,0)<\rho(x)$ at every state (\path{responsePath_gt_stage}).
Write $x^{10}_t$ and $x^{11}_t$ for the two paths. If $x^{10}_t\le x^{11}_t$,
then
$\F(x^{10}_t,\sigma(x^{10}_t,0))\le\F(x^{10}_t,\rho(x^{10}_t))\le\F(x^{11}_t,\rho(x^{11}_t))$,
because $\rho$ is nondecreasing. By induction the target-only path stays below
the forward-looking one (\path{rulePathF_le_of_rule_le_and_upper_mono}), so
$x^{10}_\infty\le x^{11}_\infty$. Equality would give $\sigma(x^{11}_\infty,0)=x^{11}_\infty=\rho(x^{11}_\infty)$, which the foresight premium
rules out, so $x^{10}_\infty<x^{11}_\infty$
(\path{targetOnly_limit_lt_dynamic_limit_of_response_monotone}; under (RM),
\path{targetOnly_limit_lt_dynamic_limit}). On the family, interiority gives
$u_{\mathrm{hi}}<q$, and the comparison reads $u_0/(q-r)<x^{11}_\infty$ from every start
(\path{targetOnly_steady_lt_dynamic_steady}). This proves (iv). \qed

\subsubsection*{I. Proof of \Cref{thm:four-cells}}

The theorem puts together results proved above and adds no new estimate. It is
stated on the enriched subclass of panel~A0 (\path{four_cell_generalF'}). Its
form with a nondecreasing response in place of (RM) is
\path{four_cell_generalF'_of_response_monotone}; \path{four_cell_generalF'} is
that form at the response of Step~4 of panel~E
(\path{responsePath_monotoneOn_of_margin}), and
\path{four_cell_generalF}, under the two-branch relief condition, follows from it
through \path{reliefSharp_of_relief}. Each case uses the following hypotheses.

\emph{Case $00$.} The myopic case, \Cref{prop:myopic-benchmark}, needs the
myopic relief condition, and here that condition follows from (RL). The sender's stage
curvature is at most its dynamic one, $\bar Q\le M_{\mathrm{tot}}$
(\path{Qbar_le_Mtot}), so
$\ell_{\mathrm{hi}}\bar Q\le\ell_{\mathrm{hi}}M_{\mathrm{tot}}<k_{\mathrm{lo}}p_{\mathrm{lo}}$
(\path{stageRelief_of_reliefSharp}). Interiority and the lifting condition enter only
through $u_{\mathrm{hi}}<q$ and $u_{\mathrm{hi}}<k_{\mathrm{lo}}$
(\path{uhi_lt_q}, \path{uhi_lt_k}). With these, Step~5 of panel~F gives a zero
realized sanction and $x_t\le(1-d_{\mathrm{lo}})^tx_0\to0$
(\path{myopicPath_tendsto_zero}).

\emph{Case $01$.} This is \Cref{prop:sender-only}(iii), Step~8 of panel~G
(\path{senderOnlyPath_tendsto_zero'}). It uses interiority, the lifting condition, (RL),
compliance neutrality, and (T6) through $M_{\mathrm{tot}}$ (Step~2 of
panel~G).

\emph{Cases $10$ and $11$.} Apply Step~5 of panel~H, with $u_{\mathrm{hi}}<q$
from interiority and with the entrenchment condition
(\path{targetOnly_limit_lt_dynamic_limit_of_response_monotone}). From each $x_0$
the target-only path converges to $x^{10}_\infty$ with $\sigma(x^{10}_\infty,0)=x^{10}_\infty$, and
the forward-looking path converges to $x^{11}_\infty$ with $\rho(x^{11}_\infty)=x^{11}_\infty$ and $x^{10}_\infty<x^{11}_\infty$. The
fixed point is positive because a non-compliant target is sanctioned at every state:
$x^{10}_\infty=\sigma(x^{10}_\infty,0)>0$ by (S11) with $u_{\mathrm{lo}}>0$
(\path{stage_argmax_pos}). Case $10$ uses only the entrenchment condition (Steps~1--3 of
panel~H). Case $11$ also uses $u_{\mathrm{hi}}<q$ and, through the nondecreasing
response, (RM). Neither case uses the lifting condition, (T6), (RL) or neutrality.

On the family (\path{PlacableEconomyFcs.four_cell_closedForm}), Step~4 of
panel~H gives $x^{10}_\infty=u_0/(q-r)$ from every start. This is positive because
$u_0>0$ and interiority gives $r<q$ (\path{steady10_pos}), and Step~5 gives a forward-looking
limit strictly above it. \qed

\subsection{Consequences for the forward-looking equilibrium of the class under the lifting condition}
\label{app:proofs:branch}

Throughout this subsection $m=\beta_S g_S d_{\mathrm{lo}}/\bar Q$ abbreviates the
bound \eqref{eq:sender-wedge} on the sender's foresight premium, $\beta_S w d_{\mathrm{lo}}/\bigl(q(1-\beta_S g_{\mathrm{lo}})\bigr)$
on the family; it is a combination of the class's rates, not a parameter. $V_S$ and $V_T$ are the values of the
constructed forward-looking equilibrium of \Cref{prop:placable-family} and
$R(x)$ its realized sanction at non-compliance.

\paragraph{The premium bound.}
The proofs below read the premium bound \eqref{eq:sender-wedge} of
\Cref{sec:results:wedge}. Its proof follows the formalization.

\begin{proof}[Proof of \eqref{eq:sender-wedge}]
Let the member satisfy interiority. Fix a value
$V\in\mathcal C_S$ in the sender's invariant set and a state $x\in[0,1]$, and write $\phi(s)=\Us(x,0,s)$,
$H(s)=\beta_SV(\F(x,s))$ and $s^\circ=\sigma(x,0)$. The response $\rho_V(x)$ at
non-compliance maximizes $\phi+H$ on $[0,1]$ (\path{response_isMaxOn}); at $V=V_S$ it is
$R(x)$. \emph{The continuation's lower slope.} For $s_1\le s_2$ in $[0,1]$ the
lower bound (T2) on the transition's sanction slope gives $\F(x,s_2)-\F(x,s_1)\ge d_{\mathrm{lo}}(s_2-s_1)\ge0$,
and every value in the invariant set rises by at least $g_S$ per unit of state
(\path{clip_lowerSlope}). Hence $H(s_2)-H(s_1)\ge L(s_2-s_1)$ with
$L=\beta_Sg_Sd_{\mathrm{lo}}$. Here $L>0$, because $\beta_S>0$, $d_{\mathrm{lo}}>0$,
and $g_S>0$ by (S2) and $\beta_Sg_{\mathrm{lo}}<1$.
\emph{The base objective.} By (S4), $\phi$ is $\bar Q$-semiconvex on $[0,1]$. It
is maximized at $s^\circ$, and $0<s^\circ<1$: $s^\circ>0$ by (S11)
(\path{stage_argmax_pos}), and $s^\circ<1$ by myopic interiority $u_{\mathrm{hi}}<q$,
which interiority implies (\path{uhi_lt_q}, \path{stage_argmax_lt_one}).
\emph{The shifted maximizer.} Every maximizer $s$ of $\phi+H$ on $[0,1]$ satisfies
$s\ge \tilde s:=\min\{1,s^\circ+L/\bar Q\}$
(\path{min_one_shifted_argmax_le_of_semiconvex}). To see this, suppose $s<\tilde s$.
If $s<s^\circ$, moving to $s^\circ$ does not lower $\phi$ and raises $H$ by at least
$L(s^\circ-s)>0$. If $s^\circ\le s<\tilde s$, write $s=\lambda s^\circ+(1-\lambda)\tilde s$ with
$\lambda=(\tilde s-s)/(\tilde s-s^\circ)$. The semiconvexity inequality at $s$ and the quadratic lower
support at the interior maximizer, $\phi(\tilde s)\ge \phi(s^\circ)-\tfrac{\bar Q}{2}(\tilde s-s^\circ)^2$
(this step uses $s^\circ>0$;
\path{quadratic_support_right_of_semiconvex_interior_max}), give
$\phi(s)\le \phi(\tilde s)+\tfrac{\bar Q}{2}(\tilde s-s)(\tilde s+s-2s^\circ)$. Because $\tilde s\le s^\circ+L/\bar Q$,
$H(\tilde s)-H(s)\ge L(\tilde s-s)\ge\bar Q(\tilde s-s^\circ)(\tilde s-s)$. Together these give
$(\phi+H)(\tilde s)-(\phi+H)(s)\ge\tfrac{\bar Q}{2}(\tilde s-s)^2>0$. Either way, $s$ is not a
maximizer.
\emph{The interior cap removes the clip.} By the cap in Step~2 of panel~E of
\Cref{app:proofs:placable}, every maximizer at non-compliance is at most $A/q$
(\path{dynamicArgmax_le_cap}, \path{responsePath_le_cap}), and $A/q<1$ by
interiority (\path{cap_lt_one}). So $\tilde s\le\rho_V(x)<1$. Hence $\tilde s=s^\circ+L/\bar Q$
and $\rho_V(x)\ge\sigma(x,0)+\beta_Sg_Sd_{\mathrm{lo}}/\bar Q$
(\path{responsePath_ge_stage_add_wedge}). Substituting
$g_S=w_{\mathrm{lo}}/(1-\beta_Sg_{\mathrm{lo}})$ gives
$\beta_Sw_{\mathrm{lo}}d_{\mathrm{lo}}/\bigl(\bar Q(1-\beta_Sg_{\mathrm{lo}})\bigr)$
(\path{wedge_eq}). On the family $\bar Q=q$ and $w_{\mathrm{lo}}=w$
(\path{toChord}), which gives the family form.
\emph{Along the paths.} Suppose in addition that the member satisfies the
entrenchment condition and that its response at full non-compliance is non-decreasing in
the state, as in \Cref{prop:steady-gap}. These hypotheses name the two
constructed equilibria. The target-only sanction rule is $\sigma(\cdot,0)$
(\Cref{prop:target-only}(i); \path{targetOnlyRule_eq_stage}) and the
forward-looking one is $R$
(\path{dynamicRule_eq_responsePath_of_response_monotone}), so
$R(x)\ge\sigma(x,0)+m$ at every state
(\path{targetOnlyRule_add_wedge_le_dynamicRule_of_response_monotone}). From a
common start, suppose $x^{10}_t\le x^{11}_t$. Then
$x^{10}_{t+1}=\F(x^{10}_t,\sigma(x^{10}_t,0))\le\F(x^{10}_t,R(x^{11}_t))
\le\F(x^{11}_t,R(x^{11}_t))=x^{11}_{t+1}$, using $\sigma(\cdot,0)\le R$,
monotonicity of $R$, (F3) and (F2). So the forward-looking state path lies above the target-only one at
every date (\path{rulePathF_le_of_rule_le_and_upper_mono},
\path{targetOnlyPath_le_dynamicPath_of_response_monotone}). The myopic sanction rule is
non-decreasing in the state by (S7) and the uniqueness in (S3)
(\path{stage_argmax_mono}), so
$\sigma(x^{10}_t,0)+m\le\sigma(x^{11}_t,0)+m\le R(x^{11}_t)$ at every date
(\path{targetOnlySanction_add_wedge_le_dynamicSanction_of_response_monotone};
the form under the increasing-differences condition is
\path{targetOnlySanction_add_wedge_le_dynamicSanction}).
\end{proof}

\begin{proposition}[The long-run gap]
\label{prop:steady-gap}
For every economy in the class satisfying interiority, every fixed point
$x^{11}_\infty\in[0,1]$ of the dynamic best response at non-compliance and every fixed point $x^{10}_\infty\le x^{11}_\infty$ of the
myopic sanction rule $x\mapsto\sigma(x,0)$ satisfy
\[
x^{11}_\infty-x^{10}_\infty\;\ge\;\frac{\beta_S\,g_S\,d_{\mathrm{lo}}}{\bar Q-r_{\mathrm{lo}}},
\]
which on the family, where $x^{10}_\infty=u_0/(q-r)$, reads
$x^{11}_\infty-u_0/(q-r)\ge qm/(q-r)=\beta_S w d_{\mathrm{lo}}/\bigl((q-r)(1-\beta_S g_{\mathrm{lo}})\bigr)$.
Under the entrenchment condition and a response at full non-compliance non-decreasing in
the state (as under the increasing-differences condition \eqref{eq:rally-margin};
\Cref{app:proofs:placable}), from every initial state the path of the
constructed forward-looking equilibrium converges to such a fixed point.
\end{proposition}

\begin{proof}
The premium bound at $x^{11}_\infty$ reads $x^{11}_\infty=R(x^{11}_\infty)\ge\sigma(x^{11}_\infty,0)+m$, and the myopic sanction rule's
slope bound (S7) with (S4) gives $\sigma(x^{11}_\infty,0)-x^{10}_\infty\ge r_{\mathrm{lo}}(x^{11}_\infty-x^{10}_\infty)/\bar Q$;
rearrange, using $r_{\mathrm{lo}}<\bar Q$: the cross difference over the unit
square, at least $r_{\mathrm{lo}}$ by (S7), is at most the upper bound on the marginal return at the top
state less the lower bound at the bottom one by (S10) and (S11), and myopic
interiority $u_{\mathrm{hi}}<q\le\bar Q$ with $u_{\mathrm{lo}}>0$ then gives
$r_{\mathrm{lo}}<\bar Q$ (\path{rlo_lt_Qbar}; \path{fixedPoint_ge_limit10_add_gap}).
On the family $x^{11}_\infty=R(x^{11}_\infty)\ge(u_0+r\,x^{11}_\infty)/q+m$; multiply by $q>0$ and divide by $q-r>0$
(\path{fixedPoint_ge_steady10_add_gap}, \path{steadyGap_eq}). The limit exists
and is a fixed point by Step~7 of the proof of \Cref{thm:regime-separation},
which reads only the monotone response
(\path{dynamicPath_limit_ge_limit10_add_gap_of_response_monotone},
\path{dynamicPath_limit_ge_steady10_add_gap_of_response_monotone}; the form under the
increasing-differences condition, \path{dynamicPath_limit_ge_steady10_add_gap}, is a corollary). The bound is a statement about
the constructed equilibrium at fixed primitives; it is not a comparative
static in $\beta_S$ of equilibria in which both governments re-optimize, and the limit may depend on the
start unless the closed loop contracts.
\end{proof}

\begin{remark}[The premium bound is attained]
\label{rem:wedge-sharp}
At the affine member with rate and both discount factors $\tfrac12$, $q=2$,
$u_0=w=1$ and $r=0$, the constructed sender value is $\tfrac43x+\tfrac89$, the
dynamic best response to non-compliance is $\tfrac23$ at every state, the stage
vertex is $\tfrac12$ and $m=\tfrac16$, so \eqref{eq:sender-wedge} holds with
equality (\path{placableBoundary_wedge_attained}); the value is pinned down by
uniqueness of the fixed point of the sender's Bellman contraction, not posited.
Hence no constant larger than one can multiply $m$ in a bound valid at every
member (\path{wedge_coefficient_cannot_be_raised}). This is an existence
statement at one affine member: it says nothing about tightness elsewhere, and
a transition supplied with looser slope bounds yields a smaller $m$ for
the same game.
\end{remark}

\begin{proposition}[The path does not depend on the target's parameters]
\label{prop:branch-invariance}
Let two games share the transition $\F$ on the square, the sender's discount
factor $\beta_S$ and the sender's stage payoff at full non-compliance,
$\Us(x,0,s)$ on the square (on the family, the coefficients $q,u_0,r,w$). Then
any full-non-compliance equilibria in $\mathcal E$ of the two have the same sender value
and the same sanction rule at every feasible state, and their state
paths and realized sanctions coincide from every initial state and at every
date, whatever the two games' target payoffs and discount factors and sender
payoffs away from full non-compliance. On economies in the class satisfying their increasing-differences
and entrenchment conditions the constructed equilibria are such equilibria.
\end{proposition}

\begin{proof}
At non-compliance the two games' sender objectives $\Us(x,0,s)+\beta_SV(\F(x,s))$
are the same function of $(V,x,s)$ on the square (on the family, because the
compliance--sanction cross term $-kas$ vanishes), so their Bellman operators at the non-compliance
policy coincide (\path{bellmanS_BCF_zero_eq_of_defiance_section_eq}). Each
equilibrium's sender value is a fixed point of its operator at its own policy,
which is full non-compliance, hence of this one
(\path{V_S_eq_V_S_zero_of_policy_zero}); the operator is a contraction of the
whole space of bounded continuous values, so the two values coincide
(\path{senderValue_eq_of_defiance_section_eq}). The sanction rules are then the
greatest maximizers of the same objective
(\path{equilibriumRule_eq_of_defiance_section_eq}), and the paths coincide by
induction (\path{equilibriumPath_eq_of_defiance_section_eq},
\path{defiance_branch_invariance_generalF}). No theorem-specific condition, invariant set or argmax
uniqueness is read, and the statement is conditional on the two equilibria;
for the constructed ones see \path{dynamicSanction_eq_of_sender_side_eq}. It
does not say that a full-non-compliance equilibrium survives an arbitrary parameter change, it
does not concern the target's value, which does differ, and off-path responses
can differ when the payoffs away from non-compliance do.
\end{proof}

\begin{proposition}[Relief participation]
\label{prop:relief-participation}
Let an economy in the class satisfy the entrenchment condition, and let the
sender's response at full non-compliance at its constructed value be non-decreasing
in the state (as under the increasing-differences condition \eqref{eq:rally-margin}). Let
$E=(\Vt,\Vs,\alpha)$ be any equilibrium in $\mathcal E$ at which the target fully
maintains non-compliance; under these hypotheses one exists, the constructed equilibrium. Write
$s_0=\min\{1,u_{\mathrm{lo}}/\bar Q\}$ and $G_S=s_0(u_{\mathrm{lo}}-\bar Qs_0/2)>0$,
the stage-gain floor ($u_0^2/(2q)$ on the family). Fix $x_0\in[0,1]$, a length $N\ge1$ and a compliance sequence $a_t\in[0,1]$, let
$y_0=x_0$, $y_{t+1}=\F(y_t,0)$ be the state under the zero sanctions of a temporary suspension of sanctions, and
for $i\in\{S,T\}$ let
\[
J_i=\sum_{t<N}\beta_i^{\,t}\,U_i(y_t,a_t,0)+\beta_i^{\,N}V_i(y_N)
\]
be government $i$'s value of the spell followed by $E$. Then
full non-compliance, $a_t\equiv0$, maximizes $J_T$ over compliance sequences, and at
it $(J_S-V_S(x_0))+(J_T-V_T(x_0))\le-G_S$, so there is no number $\tau$ with
$V_S(x_0)\le J_S-\tau$ and $V_T(x_0)\le J_T+\tau$. If moreover no compliance
level raises the sender's payoff above full non-compliance's at zero sanction
(automatic on the family), then for every compliance sequence
\[
V_S(x_0)-J_S\;\ge\;G_S\sum_{t<N}\beta_S^{\,t}\;>\;0,
\qquad
V_T(x_0)-J_T\;\ge\;\ell_{\mathrm{lo}}\sum_{t<N}\beta_T^{\,t}a_t\;\ge\;0 .
\]
and the joint loss and the absence of $\tau$ hold at every sequence.
\end{proposition}

\begin{proof}
\emph{One period, sender.} The myopic sanction $\nu=\sigma(y,0)$ is feasible, so
the Bellman inequality gives
$V_S(y)\ge\Us(y,0,\nu)+\beta_SV_S(\F(y,\nu))$. The value is non-decreasing in
the state and $\F$ in the sanction, so $V_S(\F(y,\nu))\ge V_S(\F(y,0))$; the
stage gain of $\nu$ over zero is at least $G_S$ by (S11) at $s_0$
(\path{stage_gain_ge}), and $\Us(y,a,0)\le\Us(y,0,0)$ for every $a$
(\path{sender_relief_step}). On the family $\nu=(u_0+ry)/q$ and
$\Us(y,0,\nu)-\Us(y,a,0)=(u_0+ry)^2/(2q)\ge u_0^2/(2q)$
(\path{toChord_stageGainFloor}). \emph{One period, target.} The
Bellman equality at non-compliance, with the sanction rule $R$, gives
$V_T(y)=\Ut(y,0,R(y))+\beta_TV_T(\F(y,R(y)))$, hence
\begin{align*}
V_T(y)-\Ut(y,a,0)-\beta_TV_T(\F(y,0))
&\ge\ell_{\mathrm{lo}}a-cR(y)+\beta_T\bigl[V_T(\F(y,R(y)))-V_T(\F(y,0))\bigr]\\
&\ge\ell_{\mathrm{lo}}a+\bigl(\beta_Td_{\mathrm{lo}}g_T-c\bigr)R(y)\ge\ell_{\mathrm{lo}}a,
\end{align*}
by (P1), (P2), the lower state slope $g_T$ of $V_T$, the lower bound on the transition's
sanction slope and the entrenchment condition $c\le\beta_Td_{\mathrm{lo}}g_T$
(\path{target_relief_step}). \emph{The spell.} Multiplying the one-period
inequalities by $\beta_i^{\,t}$ and summing over $t<N$ telescopes
(\path{Foundations/FiniteDiscountedTelescope.lean};
\path{senderReliefValue_add_loss_le}, \path{targetReliefValue_add_cost_le}),
and the remaining claims are arithmetic: at non-compliance the sender's step needs
no condition at zero sanction (\path{targetReliefValue_le_defiance},
\path{relief_joint_loss_at_defiance},
\path{no_date_zero_transfer_supports_relief_at_defiance}), and at every
sequence it reads the condition that no compliance level raises the sender's
payoff at zero sanction (\path{relief_joint_loss},
\path{no_date_zero_transfer_supports_relief},
\path{relief_participation_of_margin}). Every full-non-compliance equilibrium has the
constructed values, so the argument applies to each
(\path{relief_participation_of_policy_zero_at_defiance},
\path{relief_participation_of_policy_zero}); the increasing-differences condition entered only
through the non-decreasing response and to name the constructed equilibrium.
\end{proof}

The proposition excludes one arrangement: uncompensated zero sanctions for
finitely many periods, then the same full-non-compliance equilibrium. It says nothing
about a different continuation equilibrium, a third party's subsidy, payments
at later dates under unequal discounting, or population welfare, and the last
clause presupposes transfers valued one for one by both governments, which the
game does not have. The target's half rests on the entrenchment condition and
fails outside it: at the economy of \Cref{app:formalization} in which that
condition fails and non-compliance is nonetheless the unique optimum, a one-period
temporary suspension of sanctions is worth exactly $\tfrac1{18}$ more to the target than
the equilibrium, at every state
(\path{lossOffset_relief_period_benefits_target}) --- compliance is costly, but
relief that is simply imposed requires none.

\subsection{The bracket in the target's discount factor, and uniqueness in the monotone class}
\label{app:proofs:threshold}

Round AP made the lifting condition and the entrenchment condition hypotheses of the
propositions that read them; on one structure both the relief and the non-compliance
equilibria of the forward-looking target can then be stated. The affine family of
earlier versions without state--sanction cross terms (an exact threshold $D/(\delta v+(1-\delta)D)$ in the
target's discount factor) is retired: its content is the instance of the following
proposition in which the general intermediate range closes to a point.

Write, on an economy in the class with semiconvexity in the sanction (T6), $\mu$
for the strong-concavity modulus and $M_{\mathrm{tot}}$ for the semiconvexity
ceiling of the sender's dynamic objective in the sanction, $L_{\mathrm{resp}}/\mu$
for the response's state Lipschitz rate,
$L_\alpha=M_{\mathrm{tot}}L_{\mathrm{resp}}/(k_{\mathrm{lo}}\mu)$,
$K_T^{\mathrm{rel}}=(v_{\mathrm{hi}}+\ell_{\mathrm{hi}}L_\alpha)/(1-\beta_Tg_{\mathrm{hi}})$ and
$K_T=L_{U_T}/(1-\beta_TL_{\mathrm{loop}})$ (\Cref{app:closedforms}).

\begin{proposition}[The bracket]
\label{prop:bracket}
\begin{enumerate}[label=(\roman*)]
\item \emph{Non-compliance side.} At any target value with lower state slope $g\ge0$
and any value in the sender's invariant set, full non-compliance is a target optimum at every
state whenever $c-\ell_{\mathrm{lo}}\mu/k_{\mathrm{hi}}\le\beta_Td_{\mathrm{lo}}g$,
and the strict unique optimum when the inequality is strict; since
$\ell_{\mathrm{lo}}>0$, the entrenchment condition implies the strict form at
$g=g_T$. The target's value at the
non-compliance profile carries the slope $g$ whenever the sender's response at its
constructed value is nondecreasing in the state (as under the increasing-differences condition \eqref{eq:rally-margin}) and
$g(1-\beta_Tg_{\mathrm{lo}})+\max\{0,c-\beta_Td_{\mathrm{lo}}g\}L_{\mathrm{resp}}/\mu\le v_{\mathrm{lo}}+n_{\mathrm{lo}}\underline s$:
$g_T$ under the entrenchment condition, and $\min\{g_T,g_{\mathrm{eff}}\}$ with
$g_{\mathrm{eff}}=(v_{\mathrm{lo}}+n_{\mathrm{lo}}\underline s-cL_{\mathrm{resp}}/\mu)/(1-\beta_T(g_{\mathrm{lo}}+d_{\mathrm{lo}}L_{\mathrm{resp}}/\mu))$
in general, when $cL_{\mathrm{resp}}/\mu<v_{\mathrm{lo}}+n_{\mathrm{lo}}\underline s$ and
$\beta_T(g_{\mathrm{lo}}+d_{\mathrm{lo}}L_{\mathrm{resp}}/\mu)<1$, the condition
that keeps the denominator of $g_{\mathrm{eff}}$ positive.
With that non-decreasing response (under the increasing-differences condition
\eqref{eq:rally-margin} in particular) and the non-compliance side at that
slope, in its weak form, the composed non-compliance equilibrium of
\Cref{prop:placable-family} exists with those values.
\item \emph{Relief side and no-non-compliance side.} Under interiority, the lifting condition,
compliance neutrality and the relief side
$\ell_{\mathrm{hi}}M_{\mathrm{tot}}\le k_{\mathrm{lo}}(p_{\mathrm{lo}}-\beta_Td_{\mathrm{hi}}K_T^{\mathrm{rel}})$,
the zero level then being the unique target optimum when the inequality is
strict, there is
an equilibrium in $\mathcal E$ in which the forward-looking target complies at the
zero level $a^{01}(x)$ at every state, the realized sanction is zero, the
sender's value is the zero-sanction value and the target's is the relief
value, which is $K_T^{\mathrm{rel}}$-Lipschitz because the zero level is
$L_\alpha$-Lipschitz in the state. Under interiority, the lifting condition,
$\beta_TL_{\mathrm{loop}}<1$ and the no-non-compliance side
$\ell_{\mathrm{hi}}M_{\mathrm{tot}}<k_{\mathrm{lo}}(p_{\mathrm{lo}}-\beta_Td_{\mathrm{hi}}K_T)$, full non-compliance is
not a target best reply at the non-compliance-profile values, so no equilibrium with
those values has the target maintaining non-compliance.
\item \emph{An intermediate range, not a point.} Every slope that the target's value
is shown in (i) to carry is at most
$g_T$, so every sufficient condition for non-compliance implies the non-compliance side at $g_T$; the
statements below are about that side, which is itself a sufficient condition where
the value is shown in (i) to carry $g_T$: under the entrenchment condition, or where the response does
not move with the state ($L_{\mathrm{resp}}=0$). The non-compliance side at $g_T$ is
disjoint from the relief side and, under $\beta_TL_{\mathrm{loop}}<1$, from the
no-non-compliance side. Across members that differ only in $\beta_T$, the sufficient
condition for non-compliance and the non-compliance side at $g_T$ are inherited by a more patient target,
the relief side by a less patient
one, and the no-non-compliance side by a less patient one under
$\beta_TL_{\mathrm{loop}}<1$ at the more patient member; in their strict forms,
the non-compliance and relief sides are open conditions on $\beta_T$. Across members that differ only
in the target's sanction-indexed state rates and pain floor, the non-compliance side
is inherited by a larger lower envelope slope $n_{\mathrm{lo}}$, the relief side
by a larger pain floor $p_{\mathrm{lo}}$, and, under the loop condition, the
no-non-compliance side by a larger $p_{\mathrm{lo}}$ with a smaller $n_{\mathrm{hi}}$
(on the family, where $n_{\mathrm{lo}}=n_{\mathrm{hi}}=n$ and $p_{\mathrm{lo}}=c-n$,
the non-compliance side by the larger cross term $n$ and the relief and
no-non-compliance sides by the smaller). The non-compliance side at
$g_T$ depends on the sender only through $\mu$, the upper substitution rate
$k_{\mathrm{hi}}$ and $\underline s$, and the relief side only through
$M_{\mathrm{tot}}$, the lower substitution rate $k_{\mathrm{lo}}$ and
$K_T^{\mathrm{rel}}$; both are equalities of the conditions with those constants
held fixed. Under fourteen tightness conditions --- tight slope bounds
($g_{\mathrm{lo}}=g_{\mathrm{hi}}$, $d_{\mathrm{lo}}=d_{\mathrm{hi}}=\delta$), zero
moduli ($\sigma_V=M_s=\tau_x=M_{\mathrm{cx}}=M_{\mathrm{cs}}=M_{\mathrm{mix}}=0$),
a zero upper envelope slope ($n_{\mathrm{hi}}=0$; then
$r_{\mathrm{hi}}=n_{\mathrm{lo}}=0$ follow) and the tight pairs
$\bar Q=q$, $k_{\mathrm{lo}}=k_{\mathrm{hi}}=k$,
$\ell_{\mathrm{lo}}=\ell_{\mathrm{hi}}=\ell$, $p_{\mathrm{lo}}=c$,
$v_{\mathrm{lo}}=v_{\mathrm{hi}}=v$ --- the tight slope bounds pin $\F$ to the affine law
$\F(x,s)=(1-\delta)x+\delta s$ on the square (either tight pair alone does, with
the fixed diagonal), $\mu=M_{\mathrm{tot}}=q$ and
$K_T^{\mathrm{rel}}=K_T=g_T$, $L_{\mathrm{resp}}=0$ so that the value is shown in (i) to carry $g_T$,
the tight slopes sum to one, and the non-compliance side
holds exactly when $\beta_T>\beta_T^*$ and the other two exactly when
$\beta_T<\beta_T^*$, with $\beta_T^*=D/(\delta v+(1-\delta)D)$ and $D=c-\ell q/k>0$:
the exact threshold of the affine setting without cross terms is where the
intermediate range closes.
The equalities are sufficient for this closing and are not shown to be
necessary; the sender's $w$ and $u$ pairs need not be tight.
\end{enumerate}
\end{proposition}

\begin{proof}
(i) The target's deviation from non-compliance to compliance $a$ lowers the response
by at most $k_{\mathrm{hi}}a/\mu$: a unit of compliance perturbs a
$\mu$-strongly concave objective with cross rate at most $k_{\mathrm{hi}}$
(\path{response_sub_le_of_le}, through
\path{argmax_dist_le_of_strongConcave_lipschitz_perturbation_sharp}). The
pain saved is at most $c$ per unit of that fall, the continuation lost at
least $\beta_Td_{\mathrm{lo}}g$ per unit through the lower bound on the transition's sanction slope, and
the compliance costs at least $\ell_{\mathrm{lo}}a$ (\path{targetObj_strictAntiOn_of_sharp},
\path{targetArgmax_eq_zero_of_sharp}, \path{sharp_of_lossOffset}); with a weak
inequality the objective is still non-increasing in compliance, so non-compliance is
an optimum, which is all existence needs
(\path{targetObj_antitoneOn_of_sharp_le}, \path{zero_isMaxOn_joint_of_sharp}). The
target operator at a sender value in the two-sided invariant set, with a monotone
response, preserves the lower-slope cone at slope $g$ under the slope condition, the
pain term along the response priced by $L_{\mathrm{resp}}/\mu$
(\path{targetOp_mem_lowerSlopeCone}); the non-compliance-profile value is the fixed
point of that contraction on the whole value space and lies in the cone
(\path{VTdef}, \path{VTdef_lowerSlope}, Stokey--Lucas Corollary~1 to
Theorem~3.2). Under the entrenchment condition $g_T$ balances (\path{gT_balance}); in general
$\min\{g_T,g_{\mathrm{eff}}\}$ does (\path{gSharp_balance}: when $L_{\mathrm{resp}}=0$
any slope at most $g_T$; otherwise $g_T\le g_{\mathrm{eff}}$ is equivalent to
the entrenchment condition at $g_T$, and $g_{\mathrm{eff}}<g_T$ makes the pain term at
$g_{\mathrm{eff}}$ nonnegative). The equilibrium is
\path{equilibriumTop_sharp}, \path{equilibriumTop_gSharp}.
(ii) The zero level is defined at any value in the two-sided invariant set with the sharp
rate of sanction reduction below it (\path{zeroLevelAt}, \path{response_ge_of_lt_zeroLevelAt});
across states the response at a fixed compliance moves by at most
$L_{\mathrm{resp}}|\Delta x|/\mu$ (\path{response_dist_le_state}), so the zero
level is $L_\alpha$-Lipschitz (\path{zeroLevelAt_lipschitz}). The relief
operator $V\mapsto\Ut(y,a^{01}(y),0)+\beta_TV(\F(y,0))$ is a contraction
preserving the $K_T^{\mathrm{rel}}$-Lipschitz clause (\path{reliefOp_lipschitz},
\path{VTrel_lipschitz}); the zero level is the unique argmax at the relief
values (\path{zeroLevel_argmax_unique_rel}) and the equilibrium is
\path{reliefEquilibriumTop}, with \path{reliefEquilibriumTop_rule_eq_zero}; at
the weak relief side the zero level is still an optimum and the equilibrium
exists (\path{jointRel_le_of_lt}, \path{reliefBestReply_of_relief_side_le},
\path{exists_relief_equilibrium_iff}). The
non-compliance-profile value is $K_T$-Lipschitz with no further condition
(\path{VTdef_lipschitz}), and the deviation to the zero level at the
constructed sender value is profitable under the no-non-compliance side
(\path{defiance_not_maxOn_of_relief_side},
\path{no_defiance_equilibrium_of_relief_side}).
(iii) A slope satisfying the slope condition of (i) is at most $g_T$, since
the pain term is nonnegative (\path{le_gT_of_balance},
\path{sharp_gT_of_sharp_of_balance}); at $L_{\mathrm{resp}}=0$ the pain term
vanishes and $g_T$ balances (\path{gT_balance_of_Lresp_eq_zero}), a situation that
is affine by hypothesis, since a zero response slope pins the transition to
$\F(x,s)=\F(1,0)\,x+(1-\F(1,0))\,s$ on the square
(\path{F_eq_affineF_of_Lresp_eq_zero}); and a slope
that balances and satisfies the non-compliance side at $\beta_T$ does both at every
larger $\beta_T$, both sides of the balance moving the right way and the
sender's constants not depending on $\beta_T$
(\path{certified_sharp_mono_}$\beta$\path{T}), while the sender's value, and so
its response, is the same at both members (\path{VSstar_eq_of_sender_side_eq}). \path{bracket_existence_sides_disjoint} for the relief side: the lower
slope bounds give $g_{\mathrm{lo}}+d_{\mathrm{lo}}\le1$ (\path{glo_add_dlo_le_one}), (P2)
and (P3) at compliance zero give
$v_{\mathrm{lo}}+n_{\mathrm{lo}}\underline s\le v_{\mathrm{hi}}+(c-p_{\mathrm{lo}})\underline s$
(\path{vlo_nlo_floor_le}), and together they give
$\beta_Td_{\mathrm{lo}}g_T\le\beta_Td_{\mathrm{hi}}K_T^{\mathrm{rel}}+c-p_{\mathrm{lo}}$, with no
neutrality, the lifting condition or loop condition read;
\path{bracket_sides_disjoint} for the no-non-compliance side (by $\mu\le M_{\mathrm{tot}}$,
$d_{\mathrm{lo}}\le d_{\mathrm{hi}}$, $g_T\le K_T$);
\texttt{sharp\_side\_mono\_$\beta$T}, \texttt{relief\_side\_anti\_$\beta$T},
\texttt{relief\_side\_notdef\_anti\_$\beta$T}; \path{sharpSideSet_isOpen},
\path{reliefSideSet_isOpen}; \path{sharp_side_mono_nlo}, \path{relief_side_anti_plo},
\path{relief_side_notdef_anti};
\path{sharp_side_eq_of_target_side_eq}, \path{relief_side_eq_of_target_side_eq};
\path{collapse_mu}, \path{collapse_Mtot}, \path{collapse_KTrel}, \path{collapse_KT},
\path{collapse_sides}; the tight slopes sum to one
(\path{glo_add_dlo_eq_one_of_tight}), and either tight pair pins the affine law on
the square (\path{eq_of_sanctionChords_tight}, \path{eq_of_stateChordBandAll_tight},
\path{add_eq_one_of_tight}; on the class \path{F_eq_affineF_of_tight_sanction},
\path{F_eq_affineF_of_tight_state}); $L_{\mathrm{resp}}=0$ there
(\path{collapse_Lresp}), and the threshold form is
\path{collapse_threshold}.
\end{proof}

The three sides are inhabited, at the affine member and off the affine law:
the members $\beta_T=\tfrac12$ of \Cref{app:formalization} are on the non-compliance
side, and the same coefficients at $\beta_T=\tfrac18$ are on the relief and
no-non-compliance sides (\path{placableAffineLo_reliefEquilibrium},
\path{placableQuadDiagLo_reliefEquilibrium}, \path{placableQuadDiag_sharpEquilibrium},
\path{exists_bracket_both_sides_nonaffine}). An exact-arithmetic evaluation of
the sides on a grid of step $1/400$
(\path{docs/archive/workflow-scratch/bracket_probe3_2026-09-23.py}) finds, at
the affine member, the relief side up to $0.180$, the non-compliance side at $g_T$
from $0.185$ and the sufficient condition for non-compliance, at $\min\{g_T,g_{\mathrm{eff}}\}$ below the
entrenchment condition, from $0.1875$; and at the non-affine member the relief side up to
$0.145$, the non-compliance side at $g_T$ from $0.2325$ and the sufficient condition for non-compliance from
$0.235$. The intermediate range, between the relief side and the sufficient condition for
non-compliance, is open at
both. At the tight-modulus affine member without cross term ($\sigma_V=0$, $r=0$) the
response does not move with the state, so the value is shown in (i) to carry $g_T$
without the entrenchment condition
(\path{gT_balance_of_Lresp_eq_zero}), and the sides are complementary at
$\beta_T=2/11$. The intermediate range is therefore not the price of the transition alone:
it is open at an affine member, which has the sender cross rate $r=1/2$ and the
supplied modulus $\sigma_V=1$, and it closes to a point at the tight-modulus
one without cross term; in general its width depends, through the constants of
(i)--(iii), on the spreads of the slope and rate bounds, the supplied moduli,
the sender's cross rate and the target's envelope slopes. The sufficient condition for non-compliance holds on
an upper set of discount factors (\path{certified_sharp_mono_}$\beta$\path{T}).
So does the existence of a full-non-compliance equilibrium itself, under a monotone
response and a monotone non-compliance value at one member, inside the intermediate range as well
(\Cref{prop:patience-existence}). Where the intermediate range closes in (iii), the threshold is an
equilibrium threshold: a full-non-compliance equilibrium exists exactly when
$\beta_T\ge\beta_T^*$ (\Cref{prop:collapse-threshold}).
On the intermediate range these
sufficient conditions determine no equilibrium, and the equilibria of both sides are members of
$\mathcal E$, so they inherit \Cref{prop:history-deviations}. On the family, at
a member with tight moduli and no cross terms the values are affine in the state,
where the Bellman equation closes in one step \citep{Puterman1994}; off
it, and on the class wherever a rate pair is split, the values are only slope-
and Lipschitz-bounded, which is what the intermediate range
measures.

The two ends of the intermediate range can be characterized exactly in general, as thresholds in the
target's discount factor along members that differ only in it. The relief end needs one
condition on the class: the state's lower rate at zero sanction covers the compliance cost
along the zero level, $\ell_{\mathrm{hi}}L_\alpha\le v_{\mathrm{lo}}$, where
$L_\alpha=M_{\mathrm{tot}}L_{\mathrm{resp}}/(k_{\mathrm{lo}}\mu)$ is the zero level's state
Lipschitz rate.

\begin{proposition}[Exact thresholds in the target's discount factor]
\label{prop:exact-thresholds}
Consider economies in the class that differ only in the target's discount factor
$\beta_T\in(0,1)$; (ii) is stated on the relief class, and only (i) needs
semiconvexity in the sanction \textup{(T6)}.
\begin{enumerate}[label=(\roman*)]
\item With \textup{(T6)}, under interiority, the lifting condition and compliance
neutrality, an equilibrium in $\mathcal E$ whose
realized sanction is zero at every state exists if and only if the zero level is a target
best reply at the relief values. If also $\ell_{\mathrm{hi}}L_\alpha\le v_{\mathrm{lo}}$, the
discount factors at which it exists form an interval $(0,\beta_R]$ whenever they are
non-empty with supremum $\beta_R<1$.
\item If the sender's response at full non-compliance is nondecreasing in the state and the
slope condition of \Cref{prop:bracket}(i) holds at $g=0$,
$cL_{\mathrm{resp}}/\mu\le v_{\mathrm{lo}}+n_{\mathrm{lo}}\underline s$, the discount factors at
which an equilibrium in $\mathcal E$ with full non-compliance at every state exists form an interval
$[\beta_D,1)$ whenever they are non-empty with infimum $\beta_D>0$.
\item Every discount factor at which an equilibrium in $\mathcal E$ has a zero realized sanction
at every state is below $c/(d_{\mathrm{lo}}v_{\mathrm{lo}}+cg_{\mathrm{lo}})$.
\end{enumerate}
\end{proposition}

\begin{proof}
(i)~An equilibrium with a zero realized sanction at every state is the relief equilibrium:
its sender value is the zero-sanction value, its policy the zero level and its target value
the relief value (\path{relief_equilibrium_of_zero_rule}), so its policy is a best reply
there. Conversely the best-reply test builds the relief equilibrium
(\path{reliefEquilibriumOfBestReply}, \path{exists_relief_equilibrium_iff}); the relief
side of \Cref{prop:bracket}(ii) is one sufficient condition for the test. Under
$\ell_{\mathrm{hi}}L_\alpha\le v_{\mathrm{lo}}$ the stage payoff along the zero level rises with
the state, so the relief value is nondecreasing (\path{VTrel_monotone}). Along the common
zero-sanction path the relief values at $\beta_1\le\beta_2$ solve $V_i=u+\beta_iV_i\circ z$,
so $\beta_2V_2-\beta_1V_1$ is nondecreasing. A compliance below the zero level draws a sanction
and raises the successor state, so its gain is smaller at the less patient member; a
compliance above it draws none, and its gain does not depend on the discount factor. The best
reply at the more patient member therefore holds at the less patient one
(\texttt{reliefBestReply\_anti\_$\beta$T}). The best-reply inequalities are weak and the relief
values converge as the discount factor does, so the set is closed
(\path{reliefBestReply_of_tendsto}), which gives $(0,\beta_R]$
(\path{reliefExistenceSet_eq_Ioc}). (ii)~The balance at $g=0$ makes the non-compliance value
nondecreasing (\path{VTdef_monotone_of_weak_hvL}); \Cref{prop:patience-existence} makes the set
an upper set, and \Cref{prop:defiance-closed} closes it (\path{defianceExistenceSet_eq_Ico}).
(iii)~At or above that bound the strong entrenchment condition holds
(\texttt{lossOffset$_0$\_iff\_$\beta$T\_ge}) and excludes such an equilibrium
(\path{no_zero_rule_equilibrium}, \path{reliefExistence_lt_margin}). General in $\F$.
\end{proof}

The proposition does not order the two thresholds: nothing in the class excludes a relief and
a full-non-compliance equilibrium at the same discount factor. The sufficient conditions of
\Cref{prop:bracket} lie inside the two sets, the sufficient condition for relief in the first and
the sufficient condition for non-compliance, with the increasing-differences condition, in the second;
so when $\beta_R<\beta_D$ the intermediate range they leave
contains the exact intermediate range $(\beta_R,\beta_D)$.

At the affine member of the family under the lifting condition every hypothesis of \Cref{prop:exact-thresholds} holds:
$\ell_{\mathrm{hi}}L_\alpha=5/23$ against $v_{\mathrm{lo}}=30$, a response made monotone by the
increasing-differences condition, and $cL_{\mathrm{resp}}/\mu=20/23$ (\path{aw_general_hypotheses}). There the two
thresholds can be computed. Write $\beta_R=36/197\approx0.18274$ and
$\beta_D=(82-7\sqrt2)/394\approx0.18300$.

\begin{proposition}[The exact intermediate range at the affine member]
\label{prop:affine-window}
At the affine member of the family under the lifting condition, with coefficients
$(q,u_0,k,r,w,\ell,c,n,v)=(3,1,3,\tfrac12,1,1,4,0,30)$, transition
$\F(x,s)=(x+s)/2$ and $\beta_S=\tfrac12$ (\Cref{app:witness:openbox}), and target
discount factor $\beta_T\in(0,1)$:
\begin{enumerate}[label=(\roman*)]
\item an equilibrium in $\mathcal E$ whose realized sanction is zero at every state
exists if and only if $\beta_T\le\beta_R$;
\item an equilibrium in $\mathcal E$ with full non-compliance at every state exists if and
only if $\beta_T\ge\beta_D$;
\item at $\beta_1=25248/138097$, which satisfies $\beta_R<\beta_1<\beta_D$, neither
exists, and for every $t\in[0,4/9]$ there is an equilibrium in $\mathcal E$ with policy
$\alpha(x)=t+287x/4206$ and realized sanction
$\tfrac{701}{2100}\cdot\tfrac{40-90t}{29}+x/10$; distinct $t$ give distinct policies,
and for $t<4/9$ the sanction is positive at every state.
\end{enumerate}
\end{proposition}

\begin{proof}
Every value is explicit. The dynamic best response is the clipped vertex of the sender's
concave objective in the sanction (\texttt{$\sigma$\_dyn\_top\_eq\_clip}), and each
Bellman equation is checked coefficient by coefficient, so each profile is an
equilibrium of the game itself (\path{awReliefEquilibrium},
\path{awDefianceEquilibrium}, \path{awEquilibrium}). (i)~If the realized sanction is
zero at every state, the gap between the sender's value and $4x/3$ at $y$ is half the
gap at $y/2$, and it is bounded, so it vanishes (\path{eq_zero_of_eq_mul_comp}). The
target's policy is then the compliance $4/9+x/6$ at which the sanction just ends, and
its value is affine with slope $179/(6-3\beta_T)$. Maintaining non-compliance at the zero state pays
exactly when $179\beta_T/(6-3\beta_T)>6$, that is, when $\beta_T>\beta_R$
(\path{aw_zero_rule_le}, \path{aw_zero_rule_iff}). (ii)~At full non-compliance the sender's
value is the unique fixed point of its Bellman contraction,
$\tfrac{10-7\sqrt2}2x^2+(4\sqrt2-4)x+\tfrac{8(2-\sqrt2)}7$, with response
$\tfrac{8(\sqrt2-1)}7+(3-2\sqrt2)x$. A unit of compliance costs the target one unit
and, while the sanction is positive, lowers it by $3/D$, where $D=7(2+\sqrt2)/8$ is the
sender's curvature; a unit of sanction costs the target $P$, its pain less the
discounted value of the state it builds. Full non-compliance is a best reply exactly when
$P\le D/3$, and $(P-D/3)\Delta=\tfrac{197}{12}(\beta_D-\beta_T)$ with $\Delta>0$
(\path{awd_threshold}); below $\beta_D$, complying at $1/10$ at the zero state pays
(\path{aw_full_defiance_le}, \path{aw_full_defiance_iff}). So the two existence sets are
the intervals of \Cref{prop:exact-thresholds}, with endpoints $\beta_R$ and $\beta_D$, and
$\beta_R$ is below the bound of that proposition's part (iii), here $4/17$
(\path{aw_exact_thresholds}). (iii)~is \texttt{aw$\beta_1$\_mem\_window},
\path{aw_continuum} and \texttt{aw\_window\_at\_$\beta_1$}.
\end{proof}

The proposition concerns one member, with the transition $\F(x,s)=(x+s)/2$ fixed, and
is affine-only. Its first two parts are instances of \Cref{prop:exact-thresholds}. The
continuum in (iii) is not an instance of a general result: it rests on the target's
indifference among compliance levels, which needs the family's linear compliance cost and
the particular discount factor $\beta_1$; with a strictly convex cost the indifference
disappears (\Cref{prop:interior-witness}). At that member the proposition settles the
intermediate range in both directions. The
sufficient conditions are conservative there: the exact bounds lie strictly inside the intermediate range,
from about $0.180$ to $0.1875$, that they leave. And the exact intermediate range is not empty of
equilibria: at one discount factor inside it, $\beta_1$, there is a continuum of
equilibria with compliance rising in the state, while neither endpoint profile is an
equilibrium. Other policies at $\beta_1$, and the equilibria at other discount factors
inside the intermediate range, are not classified.

The same structure carries the sender-side counterpart. The lower bound
$u_{\mathrm{lo}}+\beta_Sg_Sd_{\mathrm{lo}}$ on the sender's dynamic marginal return to
sanctioning and the upper bound
$A=u_{\mathrm{hi}}+\beta_SK_Sd_{\mathrm{hi}}$ on that return bracket the lifting condition: when the lower
bound exceeds an upper substitution rate $k_{\mathrm{hi}}$ that the member
has, every equilibrium in $\mathcal E$ whose policy is
nonincreasing in the state carries a sanction of at least
$(u_{\mathrm{lo}}+\beta_Sg_Sd_{\mathrm{lo}}-k_{\mathrm{hi}})/\bar Q>0$ at every state,
full compliance included, on a member at which a larger concession is never
worth more to the sender at zero sanction
(\path{nonplacable_floor_of_lowerAppetite}, from the premium floor of
\path{equilibrium_sanction_ge_wedge_of_antitone}); when $A\le k_{\mathrm{lo}}$ full
compliance ends the sanction (\path{response_at_one}). The lower bound rises
strictly with the sender's discount factor across members that differ only in
it (\path{lowerAppetite_lt_of_}$\beta$\path{S_lt}): the floor moves with the sender's
patience. The upper bound $A$ rises strictly with it too
(\texttt{A\_lt\_of\_$\beta$S\_lt}).

\begin{proposition}[Uniqueness in the monotone class]
\label{prop:monotone-uniqueness}
On an economy in the class satisfying the strong entrenchment condition
$c\le\beta_Td_{\mathrm{lo}}v_{\mathrm{lo}}/(1-\beta_Tg_{\mathrm{lo}})$ (the entrenchment condition
at $n_{\mathrm{lo}}=0$, and stronger in general), every equilibrium in $\mathcal E$ whose
sanction rule is nondecreasing in the state has full non-compliance at every
feasible state,
the constructed sender value, the constructed target value, and the
constructed response as its rule. When the constructed response at full
non-compliance is non-decreasing in the state as well (as under the increasing-differences
condition \eqref{eq:rally-margin}), the constructed non-compliance equilibrium exists and
belongs to the class.
\end{proposition}

\begin{proof}
The target's value is
its own joint integrand at its policy (\path{V_T_eq_at_policy}). At an ordered
pair of states $y_1\le y_2$ the deficit of the value below the slope
$v_{\mathrm{lo}}/(1-\beta_Tg_{\mathrm{lo}})$ is at least $\beta_T$ times the deficit
at a successor pair (\path{slopeDeficit_step}). The comparison sets the lower state
at its own policy and sanction against the higher state's deviation to full
non-compliance, which draws a weakly higher sanction, $\Delta s\ge0$, because the rule is
nondecreasing and the dynamic best response is antitone in compliance. The value at
$y_1$ is its integrand at its policy and the value at $y_2$ is at least its
integrand at full non-compliance. The stage difference is at least
$v_{\mathrm{lo}}\Delta y-c\Delta s$: compliance at $y_1$ only costs, the pain of the
extra sanction is at most $c$ per unit, and the state is worth at least
$v_{\mathrm{lo}}$ per unit at full non-compliance, the envelope term
$n_{\mathrm{lo}}s\,\Delta y$ being nonnegative (P2, P3). The successor of $y_2$
exceeds that of $y_1$ by at least $g_{\mathrm{lo}}\Delta y+d_{\mathrm{lo}}\Delta s$,
and the strong entrenchment condition covers the pain. So the infimum $D$ of the deficit over ordered pairs, finite because
the value is bounded, satisfies $D\ge\beta_TD$ and hence $D\ge0$
(\path{target_lowerSlope_of_rule_monotone}; no operator on the value space, no
continuity of the rule and no condition on the policy is used). The deviation from $\alpha(x)$ to full
non-compliance raises the sanction, because the dynamic best response is antitone in
compliance under decreasing differences with no uniqueness
(\texttt{$\sigma$\_dyn\_top\_antitoneOn\_a\_of\_DD}), the pain of the rise is
covered by the continuation at that slope through the lower bound on the transition's sanction slope,
and the compliance saved is worth at least $\ell_{\mathrm{lo}}\alpha(x)$; so $\alpha(x)=0$
(\path{policy_eq_zero_of_rule_monotone}). The two values are then the unique
fixed points of the two contractions at non-compliance
(\path{values_eq_of_policy_zero}, \path{unique_equilibrium_of_rule_monotone}). No
compliance neutrality is read. For the last sentence, the constructed policy
is constant and its rule is the constructed response, nondecreasing by
hypothesis (under the increasing-differences condition by \path{responsePath_monotoneOn_of_margin}),
which puts the constructed equilibrium in the class
(\path{equilibriumTop_mem_monotone_class}).
\end{proof}

The class is a hypothesis on the quantifier and the proposition says nothing
outside it; on the relief side, uniqueness among all equilibria is open.
The non-compliance core holds over an arbitrary economy, not only the class:
two slope bounds on the transition, decreasing differences of the sender's payoff in
(sanction; compliance), four target rates at full non-compliance and the strong entrenchment
condition give every equilibrium in $\mathcal E$ with a nondecreasing sanction rule a
target value rising at least at rate $v_{\mathrm{lo}}/(1-\beta_Tg_{\mathrm{lo}})$ and
full non-compliance at every state, with no condition on the policy and none of the
sender's curvature, invariant-set, Riccati or interiority conditions read
(\path{lowerSlope_of_rule_monotone}, \path{policy_eq_zero_of_rule_monotone}); its
extra generality over the class is illustrated by an economy whose compliance
cost is not monotone (\Cref{prop:core-outside-chord}). Under the same condition no
equilibrium in $\mathcal E$ sanctions at zero at every state
(\path{no_zero_rule_equilibrium}). The condition holds exactly when
$\beta_T\ge c/(d_{\mathrm{lo}}v_{\mathrm{lo}}+cg_{\mathrm{lo}})$, a patience bound with
no sender parameter that is sufficient and whose converse is refuted
(\texttt{lossOffset$_0$\_iff\_$\beta$T\_ge}).

\subsection{Proofs of the patience, minmax-payoff, uniqueness and threshold propositions}
\label{app:proofs:scope}

\paragraph{Proof of \Cref{prop:patience-existence}.}
The argument runs on the relief class and reads no semiconvexity of the
sanction sections \textup{(T6)}; the enriched-subclass statements of the
headline rows are its corollaries. The two members share the transition, the
stage payoffs and the sender's discount factor, so the sender's value and its
response at full non-compliance
coincide at both (\path{VSstar_eq_of_sender_side_eq}). A full-non-compliance
equilibrium exists exactly when full non-compliance is a target best reply at the
constructed values (\texttt{exists\_defiance\_equilibrium\_iff}). Compliance
lowers the response; with the response nondecreasing in the state and the
non-compliance value nondecreasing at one member, the discounted continuation
difference that compliance gives up is nondecreasing in $\beta_T$
(\path{monotone_sub_of_solves_of_monotone_right}), so the best reply at the less
patient member is preserved at the more patient one
(\texttt{exists\_defiance\_equilibrium\_mono\_$\beta$T}; headline row 53). The
slope condition of \Cref{prop:bracket}(i), which delivers the monotone
non-compliance value, is itself inherited by the more patient member: its slope side
is antitone in $\beta_T$ (\path{balanceSide_anti},
\texttt{balance\_of\_le\_$\beta$T}, \texttt{SameBut$\beta$T.balance\_of\_balance};
headline row 57). A directional example on the quartic family shows the
comparison is not vacuous.

\paragraph{Proof of \Cref{prop:security}.}
The argument reads only the lower bound $g_{\mathrm{lo}}$ on the state slope of the
transition's zero-sanction section, the lower bound $d_{\mathrm{lo}}$ on its sanction slope
from zero, three target bounds
--- a state rate $v_{\mathrm{lo}}\ge0$ at full non-compliance and zero sanction, pain at
most $c$ per unit at full non-compliance, and compliance that does not pay against a
zero sanction (used only for the attainment in (i) and (ii)) --- and the
strong entrenchment condition; it reads nothing of the sender's objective, so it holds over any
economy with these properties (\path{Theory/Equilibrium/SecurityCore.lean}),
and the class satisfies them. Let $r_0=x_0$, $r_{t+1}=\F(r_t,0)$ and
$\underline V_T(x_0)=\sum_t\beta_T^t\Ut(r_t,0,0)$. Against any feasible pure sanction sequence
$(s_t)$, compare the state under full non-compliance at every date with the zero-sanction path: the
lower slope bounds make the excess state a discounted accumulation of the sanctions,
valued at least at rate $v_{\mathrm{lo}}$, while the sanction pain is at most $c$
per unit, so the payoff gap is at least
$(\beta_Td_{\mathrm{lo}}v_{\mathrm{lo}}/(1-\beta_Tg_{\mathrm{lo}})-c)\sum_t\beta_T^ts_t$,
which the strong entrenchment condition makes nonnegative
(\path{reliefPayoff_le_defiancePayoff}; \path{Foundations/DiscountedChordSum.lean}).
Against permanent relief, compliance only costs and moves neither the sanction
nor the state, so no compliance sequence earns more than $\underline V_T$
(\path{reliefPayoff_isLeast}). In this deterministic game with perfect
monitoring a pair of pure history-dependent plans generates an open-loop path
with the same payoff (\path{jointState_eq_openLoopPath},
\path{jointPayoff_eq_openLoopPayoff}), so an open-loop pure saddle value is the
minmax and the maxmin over pure history-dependent strategies on both sides
(\path{isLeast_ceiling_history_of_openLoop_saddle},
\path{isGreatest_floor_history_of_openLoop_saddle}); under these properties $\underline V_T$ is
both (\path{reliefPayoff_isLeast_history}, \path{reliefPayoff_isGreatest_history};
headline rows 55--56), the floor attained by full non-compliance at every date
(\path{defiance_secures_reliefPayoff_history}). Equilibrium optimality against
the deviation to full non-compliance at every date gives $\Vt(x_0)\ge\underline V_T(x_0)$ in every
equilibrium in $\mathcal E$ (\path{reliefPayoff_le_V_T}; headline row 52).

\paragraph{Proof of \Cref{prop:identification}.}
Part (i) is \Cref{prop:monotone-uniqueness} (\path{unique_equilibrium_of_rule_monotone},
\path{policy_eq_zero_of_rule_monotone}; headline rows 46 and 54). For (ii), a
sanction rule that is zero at every state is nondecreasing, so an equilibrium
with it would fall in the class of (i), which makes the constructed
response its sanction rule; that response is strictly positive at full non-compliance
(\Cref{prop:placable-family}), a contradiction (\path{no_zero_rule_equilibrium};
headline row 58).

\paragraph{Proof of \Cref{prop:collapse-threshold}.}
Five of the fourteen tightness conditions of \Cref{prop:bracket}(iii) describe
which bounds are supplied for the economy, not the game. A member of the relief class with
$\sigma_V=0$ and $d_{\mathrm{lo}}=d_{\mathrm{hi}}$ is the same game as an enriched
member with the affine law's bounds, $g_{\mathrm{lo}}=g_{\mathrm{hi}}=1-d_{\mathrm{lo}}$
and $M_s=M_{\mathrm{cx}}=M_{\mathrm{mix}}=M_{\mathrm{cs}}=0$, with the same
$\beta_T^*$, and interiority and the lifting condition carry over
(\path{collapseRecert}); so the nine tightness conditions $\sigma_V=\tau_x=n_{\mathrm{hi}}=0$,
$\bar Q=q$, $k_{\mathrm{lo}}=k_{\mathrm{hi}}$, $\ell_{\mathrm{lo}}=\ell_{\mathrm{hi}}$,
$p_{\mathrm{lo}}=c$, $v_{\mathrm{lo}}=v_{\mathrm{hi}}$ and $d_{\mathrm{lo}}=d_{\mathrm{hi}}$,
with interiority, the myopic relief condition and the lifting condition, give the threshold
on the relief class, with no \textup{(T6)}
(\path{collapse_equilibrium_threshold_le_tight}). At the re-certified member the
three sides of \Cref{prop:bracket} collapse at $\beta_T^*$ (\path{collapse_sides'};
headline row 45) and the transition is the affine law on the square
(\Cref{prop:bracket}(iii)).
Above $\beta_T^*$ the non-compliance side holds and the composed non-compliance
equilibrium exists; below it, under interiority and the lifting condition, the no-non-compliance
side holds, its loop condition $\beta_TL_{\mathrm{loop}}<1$ being automatic because
$L_{\mathrm{resp}}=0$ there, so that $L_{\mathrm{loop}}=g_{\mathrm{hi}}$
(\path{collapse_Lresp'}, \texttt{$\beta$T\_ghi\_lt\_one}), and no equilibrium in
$\mathcal E$ has full non-compliance, so the non-existence direction
ranges over every equilibrium in $\mathcal E$
(\path{collapse_equilibrium_threshold}, which excludes $\beta_T=\beta_T^*$). Members at
$\beta_T=\tfrac12$ and $\tfrac18$ lie on the two sides, with $\beta_T^*=\tfrac2{11}$
(\path{Examples/PlacableSenderCollapseWitness.lean}). At $\beta_T=\beta_T^*$ the
members just above the threshold have the equilibrium, and existence passes to the
limit (\Cref{prop:defiance-closed}; \path{exists_defiance_equilibrium_at_betaStar}),
so a full-non-compliance equilibrium exists exactly when $\beta_T^*\le\beta_T$, the
threshold attained (\path{collapse_equilibrium_threshold_le}; headline row 59).
So the discount factors with a
full-non-compliance equilibrium are exactly $[\beta_T^*,1)$
(\path{collapse_defianceExistenceSet}), and when $\beta_T^*<1$ this is the
interval $[\beta_D,1)$ of \Cref{prop:exact-thresholds}(ii) with $\beta_D=\beta_T^*$,
the general theorem applying at the re-certified member, where
$L_{\mathrm{resp}}=0$ discharges its two hypotheses
(\path{collapse_exact_thresholds}): the threshold discount factor $\beta_T^*$ is the closed form
of the exact non-compliance threshold.

\begin{proposition}[Existence of non-compliance is closed in the target's discount factor]
\label{prop:defiance-closed}
Consider economies in the class that differ only in the target's discount factor.
\begin{enumerate}[label=(\roman*)]
\item If members whose discount factors converge to $\beta_T$ each have an equilibrium in
$\mathcal E$ with full non-compliance at every state, so does the member at $\beta_T$.
\item The set of discount factors in $(0,1)$ at which such an equilibrium exists
contains its infimum whenever that infimum is positive.
\end{enumerate}
\end{proposition}

\begin{proof}
A full-non-compliance equilibrium exists exactly when full non-compliance is a target best reply at
the constructed values (\texttt{exists\_defiance\_equilibrium\_iff}), and the sender's
value and response do not depend on $\beta_T$ (\path{VSstar_eq_of_sender_side_eq}).
Along the common non-compliance path the target values solve $V_n=u+\beta_nV_n\circ z$, so
$(1-\beta_n)\|V_n-V\|\le|\beta_n-\beta_T|\,\|V\|$; every deviation gain converges, and
the weak best-reply inequalities pass to the limit
(\path{defianceBestReply_of_tendsto}, \path{exists_defiance_equilibrium_of_tendsto}).
Part (ii) applies (i) to members in the set whose discount factors decrease to the
infimum (\path{defianceExistence_sInf_mem}). General in $\F$.
\end{proof}

With \Cref{prop:patience-existence} applied at the infimum, the existence set is then
the interval from the infimum to one. Where the intermediate range closes, part~(i) applied to the members
just above $\beta_T^*$ gives existence at $\beta_T=\beta_T^*$ in \Cref{prop:collapse-threshold}.

\subsection{Pure history-dependent deviations}
\label{app:proofs:history}

\begin{proposition}[Stationary equilibria survive pure history-dependent deviations]
\label{prop:history-deviations}
\emph{(For the general theory of history-dependent plans against stationary
ones in discounted dynamic programming see \citet{Strauch1966} and
\citet{Denardo1967}; the deterministic, perfect-monitoring setting used here
reduces a history-dependent plan to a dated one.)}
Let $E=(\Vt,\Vs,\alpha)\in\mathcal E$ and write
$\sigma=\sigDynTop_{\Vs}$.
\begin{enumerate}[label=(\roman*)]
\item With the sender at $\sigma$, every pure target strategy that conditions
compliance on the whole history of states and actions yields the target, from
every initial state $x_0$, a discounted payoff of at most $\Vt(x_0)$, which
$\alpha$ attains.
\item With the target at $\alpha$, every pure sender strategy that conditions
the sanction on the whole history yields the sender at most $\Vs(x_0)$, which
$\sigma(\cdot,\alpha(\cdot))$ attains.
\item At every state $x$, after every compliance $a\in[0,1]$ on or off the
equilibrium path, for every sanction $s$ and every pure history-dependent
continuation of the sender with discounted payoff $\Pi$ from $\F(x,s)$,
\[
\Us(x,a,s)+\beta_S\,\Pi\;\le\;\Us\bigl(x,a,\sigma(x,a)\bigr)
+\beta_S\Vs\bigl(\F(x,\sigma(x,a))\bigr).
\]
\end{enumerate}
\end{proposition}

\begin{proof}
The game is deterministic with perfect monitoring. Against a fixed feedback
rule of the opponent, a pure history-dependent strategy generates from a given
start one sequence of actions; the dated plan that replays that sequence,
whatever the state, generates the same states, actions and payoff
(\path{historyDiscountedPayoff_eq_replay}, in
\path{Foundations/HistoryDependentDeviations.lean}; dated plans are themselves
history-dependent strategies, \path{historyState_toHistoryPlan}, so the class
is larger). The opponent's past actions are functions of the recorded states
and own actions, so conditioning on them adds nothing. Parts (i) and (ii) are
then the lifetime verification against dated plans of
\path{Theory/Equilibrium/TopLifetimeVerification.lean}
(\path{leader_optimal_against_history_deviation_top},
\path{follower_optimal_against_history_deviation_top}); the payoff and value
bounds that verification needs are derived from continuity on the compact
cube, not assumed. Part (iii) combines the sender's optimality at the node
$(x,a)$, which the equilibrium carries at every $(x,a)$, with part (ii) from
the successor state
(\path{follower_node_optimal_against_history_deviation_top}).
\end{proof}

The proposition holds for every element of $\mathcal E$, hence for each of the
four constructed equilibria. It says that a stationary equilibrium of the paper
is unimprovable by such deviations. It does not characterize the other
equilibria that history-dependent strategies may support, it does not cover
randomized strategies, and the target's deviation is a unilateral change of
its own actions against the sender's fixed contingent response.

\subsection{Proofs of the welfare results of \Cref{sec:welfare}}
\label{app:proofs:welfare}

The propositions of \Cref{sec:welfare}, the identity \eqref{eq:welfare-identity} and
\Cref{prop:welfare-nonidentification-four} below are machine-checked in \path{Theory/Welfare/InterestGroup.lean},
\path{Theory/Welfare/PlannerZeroSanction.lean} and
\path{Examples/ChordEconomyWelfareGeneral.lean} (\Cref{app:formalization}). None of
them reads a functional form, a sign condition of an interest-group polity, or a
parameter value. \Cref{prop:welfare-rescaling} and the identity hold for any pair of
payoff functions, \Cref{prop:zero-sanction-planner} for any bounded welfare flow
and any transition with the stated chords, and the others on the class.

\paragraph{The identity \eqref{eq:welfare-identity}.}
Substituting the two representation identities,
$\Us+\Ut=\sum_g\theta_gw_g=\sum_gw_g+\sum_g(\theta_g-1)W_g=W+\sum_g(\theta_g-1)W_g$
at every state and action (\path{payoffs_add_eq}). No sign condition and no
transition is used.

The rescaling result, which \Cref{sec:welfare:nonidentification} states in words, is the
following.

\begin{proposition}[Rescaling one group]
\label{prop:welfare-rescaling}
Take an interest-group representation and a sender group $g$, and multiply $W_g$
by a positive factor while dividing $\theta_g$ by the same factor. The governments'
payoffs, and so the game and its equilibria, are unchanged, and aggregate welfare
changes by exactly the change in $W_g$. If along two paths of states and
actions the group's welfare is higher on the first by a uniform positive margin at
every date, and aggregate welfare and $W_g$ are bounded along both, then for some
factor the planner ranks the first path strictly above the second at every
planner discount factor.
\end{proposition}

\begin{proof}[Proof of \Cref{prop:welfare-rescaling}]
Replacing $(W_g,\theta_g)$ by $(\varphi W_g,\theta_g/\varphi)$ for a factor $\varphi>0$ leaves the product
$\theta_gW_g$, and so the sender's representation identity, unchanged, and the new
weight is positive; the result is a representation of the same payoffs
(\path{rescaleS}). Aggregate welfare changes only in group $g$'s term, from $W_g$ to
$\varphi W_g$, and so becomes $W+(\varphi-1)W_g$ (\path{rescaleS_welfare}).

For the ranking, write $W^P_t,W^Q_t$ for aggregate welfare and $W^P_{g,t},W^Q_{g,t}$ for
group $g$'s welfare at date $t$ along the two paths, all at most $M$ in absolute
value, with $W^P_{g,t}\ge W^Q_{g,t}+\varepsilon$ for a margin $\varepsilon>0$. Take $\varphi=1+(2M+1)/\varepsilon$. At
every date
\[
\bigl(W^P_t+(\varphi-1)W^P_{g,t}\bigr)-\bigl(W^Q_t+(\varphi-1)W^Q_{g,t}\bigr)
\ge-2M+(\varphi-1)\varepsilon=1 .
\]
The rescaled flows are bounded, so their discounted sums converge absolutely at
every $\beta\in[0,1)$, and a strict inequality at every date between bounded flows
survives discounting, since the date-zero term has weight one, also at $\beta=0$
(\path{tsum_pow_mul_lt_of_forall_lt}, \path{exists_scale_ranks},
\path{exists_rescaleS_ranks}).
\end{proof}

\begin{proof}[Proof of \Cref{prop:welfare-nonidentification}]
\emph{The representations.} For a real $\eta$, give the sender two groups, with
welfare $2\eta s$ at weight $\tfrac12$ and $\Us-\eta s$ at weight one, and the
target two groups, with welfare $-\bigl(\Ut(x,0,s)-\Ut(x,a,s)\bigr)$ at weight two
and $\Ut(x,0,s)+\bigl(\Ut(x,0,s)-\Ut(x,a,s)\bigr)$ at weight one. The weighted sums
are $\Us$ and $\Ut$ at every state and action, and aggregate welfare is
\[
W_\eta(x,a,s)=\Us(x,a,s)+\Ut(x,0,s)+\eta s
\]
(\path{etaRep}, \path{etaRep_welfare}). Compliance is a transfer between the two
target groups.

\emph{The flows.} Along the myopic benchmark the sanction is zero at every date
(\Cref{prop:myopic-benchmark}), so its flow is
$W_\eta(x_t,a^{00}(x_t),0)=\Us(x_t,a^{00}(x_t),0)+\Ut(x_t,0,0)$ whatever $\eta$.
Along the forward-looking equilibrium of \Cref{thm:regime-separation} the target
maintains full non-compliance, and the realized sanction $s_t$ is at least $\underline s$ at every
date: the forward-looking response is at least the stage response to full
non-compliance, which is at least $\underline s$ because it is non-decreasing in the
state (\path{responsePath_ge_sfloor}). Let $M_b$ bound $|\Us(x,a,s)+\Ut(x,0,s)|$
on the cube, and let $\eta_+=(2M_b+1)/\underline s$. Writing $W^{00}_t$ and
$W^{11}_t$ for the two flows under $W_{\eta_+}$,
\[
W^{11}_t-W^{00}_t\ge-2M_b+\eta_+s_t\ge1 ;
\]
under $W_{-\eta_+}$ the same bound holds with the difference reversed. All flows
are bounded by $M_b+\eta_+$, so the strict inequalities at every date survive
discounting at every $\beta\in[0,1)$
(\path{exists_reps_reversing_myopic_forward_ranking}, with the flows
\path{repFlowMyopic} and \path{repFlowForward}). The hypotheses are those of
\Cref{prop:myopic-benchmark} and of the construction in
\Cref{thm:regime-separation}, with a non-decreasing response in place of the
increasing-differences condition; compliance neutrality, semiconvexity in the sanction and the forward-looking relief condition are not used. The two representations,
$W_{\eta_+}$ and $W_{-\eta_+}$, depend on neither the initial state nor the
planner's discount factor.
\end{proof}

\paragraph{The four cases.} The same representations reverse the welfare ranking
of all four constructed equilibria of \Cref{thm:four-cells}.

\begin{proposition}[Welfare non-identification across the four cases]
\label{prop:welfare-nonidentification-four}
Take an economy in the class that satisfies the four-case conditions of
\Cref{thm:four-cells}, with a sender response at full non-compliance at its constructed
value that is non-decreasing in the state in place of the increasing-differences
condition. There are two interest-group representations of its government payoffs,
each with two groups per country and positive constant weights, such that, at every
initial state and every common planner discount factor $\beta\in[0,1)$, discounted
aggregate welfare along the four constructed equilibria is ranked
\[
11\succ10\succ00\sim01\quad\text{by one},\qquad
00\sim01\succ10\succ11\quad\text{by the other}.
\]
\end{proposition}

\begin{proof}
Take the representations $W_\eta$ of the previous proof. The hypotheses are those
under which the four equilibria are constructed in panel~I of
\Cref{app:proofs:placable}, with a non-decreasing response in place of the increasing-differences condition. Along the equilibria of cases $00$ and $01$ the sanction
is zero at every date and the two state paths coincide
(\path{myopicPath_eq_senderOnlyPath_of_reliefSharp}); by compliance neutrality
$W_\eta(x,a,0)=\Us(x,0,0)+\Ut(x,0,0)$ does not depend on compliance, so the two
flows coincide at every date and for every $\eta$, and the two cases tie. Along the
equilibria of cases $10$ and $11$ the target maintains full non-compliance. The realized sanction
$s^{10}_t$ of case $10$ is the stage response to full non-compliance, at least
$\underline s$ because that response is non-decreasing in the state, and the
realized sanction $s^{11}_t$ of case $11$ exceeds $s^{10}_t$ by at least the foresight premium
$m$ of \eqref{eq:sender-wedge} at every date
(\path{targetOnlySanction_add_wedge_le_dynamicSanction_of_response_monotone}). Let
$\mu=\min\{\underline s,m\}>0$, let $M_b$ bound $|\Us(x,0,s)+\Ut(x,0,s)|$ on the
cube, and let $\eta_+=(2M_b+1)/\mu$. Writing $W^{ij}_t$ for the flow of case $ij$
under $W_{\eta_+}$,
\[
W^{11}_t-W^{10}_t\ge-2M_b+\eta_+\bigl(s^{11}_t-s^{10}_t\bigr)\ge1,
\qquad
W^{10}_t-W^{00}_t\ge-2M_b+\eta_+s^{10}_t\ge1 ;
\]
under $W_{-\eta_+}$ the same bounds hold with every difference reversed. All
flows are bounded by $M_b+\eta_+$, so the strict inequalities at every date
survive discounting at every $\beta\in[0,1)$
(\path{exists_reps_reversing_regime_ranking}).
\end{proof}

\begin{proof}[Proof of \Cref{prop:zero-sanction-planner}]
Let $V(x)$ be the planner's value of the path with zero sanctions and compliance
$\hat\alpha$ from $x$. It is bounded and satisfies
$V(x)=W(x,\hat\alpha(x),0)+\beta V(\F(x,0))$
(\path{discountedPayoff_stationary_eq}).

\emph{Step 1: $V$ rises with the state at a rate of at most
$K=\max\{B,0\}/(1-\beta\lambda)$.} For $x\le x'$, the two zero-sanction paths
$x_t,x'_t$ stay ordered by \textup{(F2)}, and by (iii) their gap is at most
$\lambda^t(x'-x)$ (\path{zeroPath_ordered}). By (ii),
$V(x')-V(x)=\sum_t\beta^t\bigl(W(x'_t,\hat\alpha(x'_t),0)-W(x_t,\hat\alpha(x_t),0)\bigr)
\le\sum_t\beta^t\max\{B,0\}\lambda^t(x'-x)=K(x'-x)$.

\emph{Step 2: $V$ is a Bellman supersolution of the planner's problem.} At every
feasible $(x,a,s)$, by (i), by Step~1 with $\F(x,0)\le\F(x,s)$ from \textup{(F3)},
by (iii) and since $K\ge0$,
\[
W(x,a,s)+\beta V(\F(x,s))
\le W(x,\hat\alpha(x),0)+\beta V(\F(x,0))+\bigl(\eta+\beta KL\bigr)s
\le V(x),
\]
the last step by the condition on $\eta$, with equality at $a=\hat\alpha(x)$ and
$s=0$.

\emph{Step 3: verification.} Summing the supersolution inequality along any
time-dependent plan over the first $N$ dates bounds that plan's discounted welfare
up to $N$ plus $\beta^NV(x_N)$ by $V(x_0)$; the terminal term vanishes because $V$
is bounded, and the flows are bounded, so every plan's value is at most $V(x_0)$,
which the zero-sanction path attains
(\path{stationary_optimal_against_markov_deviation_of_global_bounds},
\path{zeroSanction_planner_optimal}). Its counterpart in the Kaempfer--Lowenberg
polity, where welfare is linear in the state, is
\Cref{prop:planner-optimum-lower-bound}(i); there the case $B\le0$ needs no slope
bound on the transition.
\end{proof}

\begin{proof}[Proof of \Cref{prop:planner-efficient-persistence}]
\emph{The representation.} Give the sender two groups, with welfare $s$ and
$\Us-s$, both at weight one, and the target two groups, with welfare $\Ut$ and
$-\Ut$, at weights two and one. The weighted sums are $\Us$ and $\Ut$, and
aggregate welfare is $W=\Us$ (\path{senderWelfareRep},
\path{senderWelfareRep_welfare}).

\emph{Compliance never raises the sender's payoff.} On the class compliance lowers
the sender's marginal return to sanctioning, so
$\Us(x,a,s)-\Us(x,0,s)\le\Us(x,a,0)-\Us(x,0,0)$, and the right side is zero by
compliance neutrality: $\Us(x,a,s)\le\Us(x,0,s)$ (\path{U_S_le_defiance}).

\emph{Verification.} Let $V_S$ be the sender's constructed value against full
non-compliance and $\rho$ its dynamic best response. At every feasible $(x,a,s)$,
\[
W(x,a,s)+\beta_SV_S(\F(x,s))\le\Us(x,0,s)+\beta_SV_S(\F(x,s))\le V_S(x),
\]
with equality at $a=0$ and $s=\rho(x)$. The flows and $V_S$ are bounded, so the
verification step of the previous proof makes the stationary plan $(0,\rho)$
socially optimal at $\beta_S$ against every time-dependent plan over both actions
(\path{exists_planner_efficient_rep}). Every equilibrium in $\mathcal E$ at which
the target maintains full non-compliance has the constructed values and plays $\rho$ at every
feasible state (\path{values_eq_of_policy_zero}), so that plan is the
equilibrium's play (\path{planner_efficient_rep_at_policy_zero}). On the
quadratic--linear family at $r=n=0$ the same argument, with a linear
microfoundation, gives \Cref{cor:planner-efficient-boundary}.
\end{proof}
